%%% The main file. It contains definitions of basic parameters and includes all other parts.

%% Settings for single-side (simplex) printing
% Margins: left 40mm, right 25mm, top and bottom 25mm
% (but beware, LaTeX adds 1in implicitly)
\documentclass[12pt,a4paper]{report}
\setlength\textwidth{145mm}
\setlength\textheight{247mm}
\setlength\oddsidemargin{15mm}
\setlength\evensidemargin{15mm}
\setlength\topmargin{0mm}
\setlength\headsep{0mm}
\setlength\headheight{0mm}
% \openright makes the following text appear on a right-hand page
\let\openright=\clearpage

%% Settings for two-sided (duplex) printing
% \documentclass[12pt,a4paper,twoside,openright]{report}
% \setlength\textwidth{145mm}
% \setlength\textheight{247mm}
% \setlength\oddsidemargin{14.2mm}
% \setlength\evensidemargin{0mm}
% \setlength\topmargin{0mm}
% \setlength\headsep{0mm}
% \setlength\headheight{0mm}
% \let\openright=\cleardoublepage

%% Generate PDF/A-2u
\usepackage[a-2u]{pdfx}

%% Character encoding: usually latin2, cp1250 or utf8:
\usepackage[utf8]{inputenc}
\usepackage[T1]{fontenc}

%% Prefer Latin Modern fonts
\usepackage{lmodern}

%% Further useful packages (included in most LaTeX distributions)
\usepackage{amsmath}        % extensions for typesetting of math
\usepackage{amsfonts}       % math fonts
\usepackage{amsthm}         % theorems, definitions, etc.
\usepackage{bbding}         % various symbols (squares, asterisks, scissors, ...)
\usepackage{bm}             % boldface symbols (\bm)
\usepackage{graphicx}       % embedding of pictures
\usepackage{fancyvrb}       % improved verbatim environment
\usepackage{natbib}         % citation style AUTHOR (YEAR), or AUTHOR [NUMBER]
\usepackage[nottoc]{tocbibind} % makes sure that bibliography and the lists
			    % of figures/tables are included in the table
			    % of contents
\usepackage{dcolumn}        % improved alignment of table columns
\usepackage{booktabs}       % improved horizontal lines in tables
\usepackage{paralist}       % improved enumerate and itemize
\usepackage{xcolor}         % typesetting in color
\usepackage{mathrsfs}		% mathscr
\usepackage[h]{esvect}		% vv
\usepackage{makecell}		% makecell
\usepackage{calc}			% widthof
\usepackage{tensor}			% writing tensors with ordered indices 
%%% Basic information on the thesis
% Thesis title in English (exactly as in the formal assignment)

\usepackage{todonotes}		% for todo notes
\usepackage{caption}
\usepackage[english]{varioref} % More descriptive referencing


% Experimentation with different fonts
%\usepackage{libertinus} %\usepackage[expert]{mathdesign}

%%% Basic information on the thesis

% Thesis title in English (exactly as in the formal assignment)
\def\ThesisTitle{The connection between continuum mechanics and Riemannian geometry}

% Author of the thesis
\def\ThesisAuthor{Miroslav Burýšek}

% Year when the thesis is submitted
\def\YearSubmitted{2022}

% Name of the department or institute, where the work was officially assigned
% (according to the Organizational Structure of MFF UK in English,
% or a full name of a department outside MFF)
\def\Department{Mathematical Institute, Charles University, Prague}

% Is it a department (katedra), or an institute (ústav)?
\def\DeptType{Institute}

% Thesis supervisor: name, surname and titles
\def\Supervisor{RNDr. Michal Pavelka, Ph.D.}

% Supervisor's department (again according to Organizational structure of MFF)
\def\SupervisorsDepartment{Mathematical Institute, Charles University, Prague}

% Study programme and specialization
\def\StudyProgramme{Physics}
\def\StudyBranch{FOF}

% An optional dedication: you can thank whomever you wish (your supervisor,
% consultant, a person who lent the software, etc.)
\def\Dedication{%
Dedication.
}

% Abstract (recommended length around 80-200 words; this is not a copy of your thesis assignment!)
\def\Abstract{%
Abstract.
}

% 3 to 5 keywords (recommended), each enclosed in curly braces
\def\Keywords{%
{hydrodynamics}, {geometry}
}

%% The hyperref package for clickable links in PDF and also for storing
%% metadata to PDF (including the table of contents).
%% Most settings are pre-set by the pdfx package.
\hypersetup{unicode}
\hypersetup{breaklinks=true}
% Definitions of macros (see description inside)

%%% This file contains definitions of various useful macros and environments %%%
%%% Please add more macros here instead of cluttering other files with them. %%%

%%% Minor tweaks of style

% These macros employ a little dirty trick to convince LaTeX to typeset
% chapter headings sanely, without lots of empty space above them.
% Feel free to ignore.
\makeatletter
\def\@makechapterhead#1{
  {\parindent \z@ \raggedright \normalfont
   \Huge\bfseries \thechapter. #1
   \par\nobreak
   \vskip 20\p@
}}
\def\@makeschapterhead#1{
  {\parindent \z@ \raggedright \normalfont
   \Huge\bfseries #1
   \par\nobreak
   \vskip 20\p@
}}
\makeatother

% This macro defines a chapter, which is not numbered, but is included
% in the table of contents.
\def\chapwithtoc#1{
\chapter*{#1}
\addcontentsline{toc}{chapter}{#1}
}

% Draw black "slugs" whenever a line overflows, so that we can spot it easily.
\overfullrule=1mm

%%% Macros for definitions, theorems, claims, examples, ... (requires amsthm package)

\theoremstyle{plain}
\newtheorem{theorem}{Theorem}
\newtheorem{lemma}[theorem]{Lemma}
\newtheorem{claim}[theorem]{Claim}
\newtheorem{definition}[theorem]{Definition}
\newtheorem{proposition}[theorem]{Proposition}
\newtheorem{example}[theorem]{Example}
\newtheorem{corollary}[theorem]{Corollary}

\theoremstyle{remark}
\newtheorem*{remark}{Remark}


%%% An environment for proofs

\newenvironment{myproof}{
  \par\medskip\noindent
  \textit{Proof}.
}{
\newline
\rightline{$\qedsymbol$}
}

%%% An environment for typesetting of program code and input/output
%%% of programs. (Requires the fancyvrb package -- fancy verbatim.)

\DefineVerbatimEnvironment{code}{Verbatim}{fontsize=\small, frame=single}

%%% The field of all real and natural numbers
\newcommand{\C}{\mathbb{C}}
\newcommand{\K}{\mathbb{K}}
\newcommand{\R}{\mathbb{R}}
\newcommand{\Z}{\mathbb{Z}}
\newcommand{\N}{\mathbb{N}}

%%% Useful operators for statistics and probability
\DeclareMathOperator{\pr}{\textsf{P}}
\DeclareMathOperator{\E}{\textsf{E}\,}
\DeclareMathOperator{\var}{\textrm{var}}
\DeclareMathOperator{\sd}{\textrm{sd}}

%%% Transposition of a vector/matrix
\newcommand{\T}[1]{#1^\top}

%%% Various math goodies
\newcommand{\goto}{\rightarrow}
\newcommand{\gotop}{\stackrel{P}{\longrightarrow}}
\newcommand{\maon}[1]{o(n^{#1})}
\newcommand{\abs}[1]{\left|{#1}\right|}
\newcommand{\isqr}[1]{\frac{1}{\sqrt{#1}}}

%%% Various table goodies
\newcommand{\pulrad}[1]{\raisebox{1.5ex}[0pt]{#1}}
\newcommand{\mc}[1]{\multicolumn{1}{c}{#1}}


%%% Custom commands
\renewcommand{\d}[1]{\, \mathrm{d}#1 \,}
\renewcommand{\i}{\mathrm{i}}
\newcommand{\Cont}{C}
\newcommand{\Hilb}{\mathcal{H}}
\newcommand{\Schwa}{\mathcal{S}}
\newcommand{\Domain}{\mathrm{D}}
\newcommand{\Fourier}{\mathcal{F}}
\newcommand{\Nbhood}{\mathcal{U}}
\newcommand{\Hf}{\mathscr{H}}
\newcommand{\Af}{\mathscr{A}}
\newcommand{\ve}{\varepsilon}
\newcommand{\hypF}{\vphantom{\mathrm{F}}_1\hspace{-0.7pt}\mathrm{F}_{\!1}}

\newcommand{\Sp}{\sigma}
\newcommand{\SpP}{\sigma_{\mathrm{p}}}
\newcommand{\SpAc}{\sigma_{\mathrm{ac}}}
\newcommand{\SpSc}{\sigma_{\mathrm{sc}}}
\newcommand{\SpD}{\sigma_{\mathrm{disc}}}
\newcommand{\SpE}{\sigma_{\mathrm{ess}}}


\def\ph{\phantom}
\def\vph{\vphantom}
\def\hph{\hphantom}
\def\rzw{\mathrlap}
\def\lzw{\mathllap}
\def\czw{\mathclap}
\newcommand*{\mask}[2]{\mathord{\makebox[\widthof{\(#1\)}]{\(#2\)}}}
\newcommand{\scalemath}[2]{\scalebox{#1}{\mbox{\ensuremath{\displaystyle #2}}}}


\newcommand{\dd}[2]{\frac{\mathrm{d} #1}{\mathrm{d} #2}}
\newcommand{\pd}[2]{\frac{\partial  #1}{\partial  #2}}
\newcommand{\e}[1]{\mathrm{e}^{#1}}

\newcommand{\coloneqq}{:=}
\newcommand{\eqqcolon}{=:}

\newcommand{\digamma}{\Psi} % haha dobrý

\newcommand\numberthis{\addtocounter{equation}{1}\tag{\theequation}}

%%%%%%%%%%%%%%%%%%%%%%%%%%%%%%%%%%%%%%%%%%%%%%%%%%%%%%%%%%%%%%%
% Mejroslav

%%% My own symbols
\newcommand{\vc}[1]{\boldsymbol{#1}} % vektor
\newcommand{\mat}[1]{\mathbf{#1}} % matice

\newcommand{\norm}[1]{\left \Vert #1 \right \Vert} % norma vektoru
\newcommand{\set}[1]{ \left \lbrace #1 \right \rbrace} % množina
\newcommand{\const}{\mathrm{const}} % konstanta

%Značení derivací a integrálů
\newcommand{\der}[2]{\frac{\mathrm{d}#1}{\mathrm{d}#2}} % obyčejná derivace
\newcommand{\pder}[2]{\frac{\partial #1}{\partial #2}} % parciální derivace
\newcommand{\ppder}[3]{\frac{\partial^2 #1}{\partial #2 \partial #3}} % parciální derivace
\newcommand{\fder}[2]{\frac{\delta #1}{\delta #2}} % funkcionální derivace
\newcommand{\ffder}[3]{\frac{\delta^2 #1}{\delta #2 \delta #3}} % funkcionální derivace

\newcommand{\D}{\mathrm{d} } % integrační znamení
\newcommand{\intR}{\int_{-\infty}^{\infty}} % integrál přes reálnou osu

% Značení posloupností, limit a sum
\newcommand{\sequence}[2]{ \left \lbrace #1 \right \rbrace_{#2=1}^\infty} % posloupnost
\newcommand{\sumnorm}[1]{\sum_{#1}^\infty} 
\newcommand{\limplus}[1]{\lim_{#1 \rightarrow + \infty}}
\newcommand{\limminus}[1]{\lim_{#1 \rightarrow - \infty}}


%Značení distribucí
\newcommand{\dual}[2]{\left \langle #1 ,#2 \right \rangle} % dualita
\newcommand{\test}{\mathcal{D}} % prostor testovacích funkcí
\newcommand{\dis}{\mathcal{D'}} % prostor distribucí

\newcommand{\weakstar}{\rightharpoonup^*} % slabá* konvergence 
\newcommand{\supp}{\mathrm{supp}\,} % nosič funkcí

\hyphenation{sym-met-ry}
\hyphenation{con-ju-gate}
\hyphenation{Di-rich-let}
\hyphenation{Hamil-to-nian}
\hyphenation{Hil-bert}
\hyphenation{La-place}
\hyphenation{La-placian}
\hyphenation{Le-bes-gue}
\hyphenation{non-di-men-sion-al-i-za-tion}


% Title page and various mandatory informational pages


% TESTING
\usepackage[inline]{showlabels}  \showlabels[\tiny]{}  % Uncomment to output the content of \label commands to the document where they are used


\begin{document}

\listoftodos

%%% Title page of the thesis and other mandatory pages

%%% Title page of the thesis

\pagestyle{empty}
\hypersetup{pageanchor=false}
\begin{center}

\centerline{\mbox{\includegraphics[width=166mm]{../img/logo-en.pdf}}}

\vspace{-8mm}
\vfill

{\bf\Large BACHELOR THESIS}

\vfill

{\LARGE\ThesisAuthor}

\vspace{15mm}

{\LARGE\bfseries\ThesisTitle}

\vfill

\Department

\vfill

{
\centerline{\vbox{\halign{\hbox to 0.45\hsize{\hfil #}&\hskip 0.5em\parbox[t]{0.45\hsize}{\raggedright #}\cr
Supervisor of the bachelor thesis:&\Supervisor \cr
\noalign{\vspace{2mm}}
Study programme:&\StudyProgramme \cr
\noalign{\vspace{2mm}}
Study branch:&\StudyBranch \cr
}}}}

\vfill

% Zde doplňte rok
Prague \YearSubmitted

\end{center}

\newpage

%%% Here should be a bound sheet included -- a signed copy of the "bachelor
%%% thesis assignment". This assignment is NOT a part of the electronic
%%% version of the thesis. DO NOT SCAN.

%%% A page with a solemn declaration to the bachelor thesis

\openright
\hypersetup{pageanchor=true}
\pagestyle{plain}
\pagenumbering{roman}
\vglue 0pt plus 1fill

\noindent
I declare that I carried out this bachelor thesis independently, and only with the cited
sources, literature and other professional sources. It has not been used to obtain another
or the same degree.

\medskip\noindent
I understand that my work relates to the rights and obligations under the Act No.~121/2000 Sb.,
the Copyright Act, as amended, in particular the fact that the Charles
University has the right to conclude a license agreement on the use of this
work as a school work pursuant to Section 60 subsection 1 of the Copyright~Act.

\vspace{10mm}

\hbox{\hbox to 0.5\hsize{%
In \hbox to 6em{\dotfill} date \hbox to 6em{\dotfill}
\hss}\hbox to 0.5\hsize{\dotfill\quad}}
\smallskip
\hbox{\hbox to 0.5\hsize{}\hbox to 0.5\hsize{\hfil Author's signature\hfil}}

\vspace{20mm}
\newpage

%%% Dedication

\openright

\noindent
\Dedication

\newpage

%%% Mandatory information page of the thesis

\openright

\vbox to 0.5\vsize{
\setlength\parindent{0mm}
\setlength\parskip{5mm}

Title:
\ThesisTitle

Author:
\ThesisAuthor

\DeptType:
\Department

Supervisor:
\Supervisor, \SupervisorsDepartment

Abstract:
\Abstract

Keywords:
\Keywords

\vss}

\newpage

\openright
\pagestyle{plain}
\pagenumbering{arabic}
\setcounter{page}{1}



%%% A page with automatically generated table of contents of the bachelor thesis

\tableofcontents

%%% Each chapter is kept in a separate file
%\chapter*{List of symbols}
\todo[inline]{List of symbols.}
\addcontentsline{toc}{chapter}{List of symbols}
\noindent\vspace{-2\baselineskip}
{\renewcommand{\arraystretch}{1.3}
\begin{table}[h!]
    \begin{tabular}{c|p{0.7\linewidth}}
    \end{tabular}
\end{table}
}

\chapter*{Introduction}
\addcontentsline{toc}{chapter}{Introduction}

Geometrical formalism plays essential part in the modern classical mechanics and there are several attempts to look at other fields of physics in geometrical way, especially at hydrodynamics and thermodynamics. The Hamiltonian formalism and Poisson brackets of infinite-dimensional dynamical systems have taken hold in the mathematics community since the early seventies.

In this thesis we examine the systems of hydrodynamic type which appear in a wide range of applications, not only in hydrodynamics, but also in chemical kinetics, gas dynamics and general relativity \cite{Savoldi}. These systems are quasi-linear partial differential equations of first-order
\begin{align}
    \pder{u^i}{t} = f^{i \alpha}_j(u) \pder{u^j}{x^\alpha} \:, \quad u^i=u^i(t, x^1, \cdots, x^d) \:, \quad i = 1, \cdots, N, \quad \alpha = 1, \cdots, d \:,
\end{align}
where the standard summation rule over repeated indices is assumed. On the space of fields $u^1, \cdots, u^N$ one can build a Poisson structure via the so called bracket of hydrodynamic type, or a Dubrovin-Novikov bracket of the form
\begin{align}
    \set{u^i(x),u^j(y)} = g^{ij \alpha}(u(x)) \partial_\alpha \delta(x-y) + b^{ij \alpha}_k (u(x)) \pder{u^k}{x^\alpha} \delta(x-y) \:,
\end{align}
and more generally for arbitrary pair of functionals $J_1 = \int j_1(u,u_x,u_{xx},\cdots) \, \D x$, $J_2 = \int j_2(u,u_x,u_{xx},\cdots) \, \D x$
\begin{align}
    \set{J_1, J_2} = \int \fder{I_1}{u^i(x)} P^{ij} \fder{I_2}{u^j(x)} \, \D x \:, \quad P^{ij} = g^{ij \alpha}(u(x)) \der{}{x^\alpha} + b^{ij \alpha}_k(u(x)) \pder{u^k}{x^\alpha} \:,
\end{align}
where the first-order differential operator $P^{ij}$ is called a Hamiltonian operator or a Poisson bivector of hydrodynamic type. 


The one-dimensional case is well explored. It turns out that there is a (one would say even magical) connection between the Poisson and pseudo-Riemannian geometry. The crucial is the theorem by B. A. Dubrovin and S. P. Novikov which can be found in \cite{Dubrovin-Novikov}. The potential benefit of this thesis is to provide the proof of this theorem in detailed way due to the large ammount of technical calculations which the author could not find elsewhere. The theorem shows that the functions $g^{ij}$ and $b^{ij}_k$ form a metric and Levi-Civita connection with zero curvature, respectively. This enables writing one-dimensional systems of hydrodynamic type as
\begin{align}
    \pder{u^i}{t} = g ^{i s} (\nabla_s \nabla_j h ) \pder{u^j}{x} \:,
\end{align}
where $h$ represents the hamiltonian of the system. This forms an alternative proof of hyperbolicity, which was done by M. Pavelka, I. Peshkov and V. Klika in 2020 \cite{Pavelka}. Hyperbolicity of evolutionary equations provides local existence of solutions and also the numerical resolution which guaranties that the small perturbations of the solution due to the discretization errors do not amplify exponentially in time.

The multi-dimensional case is much more complicated and up to these days it is far from beeing completely solved.

The organisation of this thesis is the following. Chapters 1 and 2 are very consise and remind some basic facts about hydrodynamical systems and geometry. In Chapter 3 one-dimensional case is handled and the proof of the Dubrovin-Novikov theorem is done. In Chapter 4 we investigate the differences in the multi-dimensional case and end the discussion with the Mokhov classification theorem.




\chapter{Appendix}


In this chapter, we shall revise some basic definitions and propositions of the functional calculus which we will need for the proper formulation of the hydrodynamic geometry. 

Over the last century, modern mathematics has worked to develop the theory of generalized functions - distributions, including Dirac's delta distributions and operations with them. For the reader who seeks a nice introduction to this theory accompanied with several physical applications, we would recommend the book by R.P.Kanwal \cite{Kanwal}. The same applies to functional variations and the conditions under which they exist. A book written with a deep pedagogical attention on this topic is e.g. the university text by von Brunt \cite{Brunt}.
Physicists, however, operate with delta distributions as they would with "ordinary" functions, using dualities as integrals and functional variations as ordinary derivatives.


After introducing the concept of Gateaux and Frechet derivatives which are suitable for arbitrary normed space we shall take a closer look on the space of smooth functions and functionals on it.

\section{Functional derivatives}

\todo[inline]{přidat nějaký kontext}

\begin{definition}[Gateaux and Frechet derivative]
    Let $X$ be a normed space, $A: X \to \R$ a functional and $x \in \Dom(A)$.
    \begin{enumerate}
        \item Let $h \in D(A)$. Gateaux derivative $\delta A(x,h)$ in the direction $h$ is defined as
        \begin{align}
            \delta A (x,h) \coloneqq \der{}{t} A(x+th) |_{t=0} = \lim_{t \rightarrow 0} \frac{A(x+th)-A(x)}{t} \:.
        \end{align}

        \item Fréchet diferential $\DD_x A$ is defined as the unique linear bounded functional satisfying the relation
        \begin{align}
            \lim_{\norm{h} \rightarrow 0} \frac{A(x+h)-A(x)- \DD_x A(h)}{\norm{h}} = 0 \:.
        \end{align}
    \end{enumerate}
\end{definition}

We are intersted in such functionals, where $X$ is the space of smooth functions with the appropriate boundary conditions. The smoothness is required because the functionals would operate with derivatives, and the values beeing zero on the boundary is required because we shall use integration by parts.

One of possible choices is to take
\begin{align}
    X^k \coloneqq \set{f \in C^k(\Omega) \,|\, \DD^\alpha f = 0 \text{ on } \partial \Omega \text{ for } |\alpha| = 1, \cdots, k-1 } \:.
\end{align}
This space can be equipped with the usual supremum norm
\begin{align}
    \norm{f}_k \coloneqq \sum_{|\alpha| \leq k} \sup_{x \in \Omega} |\DD^\alpha f(x)| \:.
\end{align}
There are two main types of functionals on these spaces -  the integral functionals and Dirac's delta distributions.

\subsection{Integral functionals}

The first type is the integral functional
\begin{align}
    A[u] = \int_\Omega a(x,u(x),u'(x),\cdots,u^{(k)}(x)) \, \D^n x \:, \label{eq:integral-functional}
\end{align}
where $a : \R \times \R^{k} \to \R$ is an analytic function.
For simplicity we often write this using a physicist's notation
\begin{align}
    A[u] = \int a(u) \, \D^n x \:,
\end{align}
meaning that $a(u)$ is in fact function $a(x,u(x),u'(x),\cdots, u^{(k)}(x))$ and also the domain $\Omega$ is specified by the context. We call $a(u)$ the density of functional $A[u]$.

Gateaux derivative of $A$ is computed as follows:
\begin{align}
    \delta A[u,h] = \int_\Omega \sum_{i=1}^k \pder{a}{u^{(i)}}(x) h^{(i)}(x) \, \D^n x \:. \label{eq:Gateaux}
\end{align}
The functional $\delta A[u,h]$ is linear and bounded in $h$, so the Fréchet derivative of $A$ exists and is given by
\begin{align}
    \DD_u A [h] = \delta A[u,h] \:.
\end{align}
To eliminate the derivatives of $h(x)$ in \eqref{eq:Gateaux} we use integration by parts in each term of the sum (exploiting the fact that boundary terms are zero), obtaining
\begin{align}
    \DD_u A [h] = \int_\Omega \left[ \pder{a}{u}(x) - \der{}{x} \pder{a}{u'}(x) + \cdots + (-1)^k \der{^k}{x^k} \pder{a}{u^{(k)}}(x) \right] h(x) \, \D^n x \:. \label{eq:perpartes}
\end{align}

It is natural to give a special name to the integrand in \eqref{eq:perpartes}.

\begin{definition}[Functional derivative of the integral functional]
    Let $A$ be the integral functional given by \eqref{eq:integral-functional}. The functional derivative of $A$ is the function
    \begin{align}
        \fder{A}{u(x)} = \pder{a}{u}(x) - \der{}{x} \pder{a}{u'}(x) + \der{^2}{x^2} \pder{a}{u''} - \cdots + (-1)^k \der{^k}{x^k} \pder{a}{u^{(k)}}(x) \:.
    \end{align}
\end{definition}

The functional derivative can be seen as the density of the integral functional $\DD_u A$ applied on $h$: 
\begin{align}
    \DD_u A[h] = \int_\Omega \fder{A}{u(x)} h(x) \, \D x \:.
\end{align}
Since the Fréchet derivative is always linear bounded functional, another way to look at this relation is the following: $\DD_u A$ is a regular distribution represented by some smooth function $\fder{A}{u(x)}$, i.e.
\begin{align}
    \DD_u A[h] = \dual{T_{\fder{A}{u(x)}}}{h(x)} \:.
\end{align}

It follows from theory \cite{Brunt} that the functional derivative beeing zero is a necessery condition for the function beeing the extremal point of the functional (given that the corresponding function is smooth enough).

\begin{theorem}[Euler-Lagrange equations]
    Let $v \in C^\infty(\Omega)$ such that $v(x)$ is stationary point of the functional $A$ in \eqref{eq:integral-functional}, i.e
    \begin{align}
        \delta A[v,h] = 0 \:.
    \end{align}
    Then  
    \begin{align}
        \fder{A}{v(x)} = 0 \:.
    \end{align}
\end{theorem}

The computation above gives meaning to the 
\begin{align}
    A[u+h] = A[u] + \dual{\fder{A}{u}}{h} + o(h)\:,
\end{align}
where
\begin{align}
    \lim_{\norm{h} \rightarrow 0} \frac{o(h)}{\norm{h}} = 0 \:.
\end{align}
In physicist's notation, the function $h$ is often denoted as $\delta u$.


\subsection{Distributions}

\begin{definition}[Distributions]
    Let $\test(\Omega)$ be the space of smooth functions with compact support in $\Omega \subseteq \R^n$. The space of distributions $\dis(\Omega)$ is the dual space, i.e. the space of bounded linear functionals on $\test(\Omega)$. An element $T \in \dis(\Omega)$ is called a distribution.
\end{definition}

\begin{definition}[Regular distribution]
    Let $a \in L^1_{\text{loc}}(\Omega)$. A regular distribution is a distribution $T_a \in \dis(\Omega)$ defined as
    \begin{align}
        \dual{T_a}{\phi(x)} = \int_\Omega a(x) \phi(x) \D^n x \:.
    \end{align}
\end{definition}

Note that if two regular distributions satisfy $T_a = T_b$, then $a=b$ almost everywhere and vice versa.

\begin{definition}[Dirac delta distribution]
    Let $y \in \Omega$. The Dirac delta distribution $\delta(x-y)$ is defined as
    \begin{align}
        \dual{\delta(x-y)}{\phi(x)} = \phi(y) \:. 
    \end{align}
\end{definition}

\begin{definition}[Derivatives of distributions]
    Let $T \in \dis(\Omega)$ be a distribution. We define its derivative by
    \begin{align}
        \dual{D^\alpha T}{\phi} = (-1)^{|\alpha|} \dual{T}{D^\alpha \phi} \:,
    \end{align}
    where $\alpha$ is a multiindex and
    \begin{align}
        D^\alpha f = \frac{\partial^{|\alpha|} f}{\partial x_1^{\alpha_1} \partial x_2^{\alpha_2} \cdots \partial x^{\alpha_N}_N} \:.
    \end{align}
\end{definition}

The reason, why distributions have such a "nice" behaviour and are often used by mathematicians, is that every distribution is diferentiable infinitely many times. However, if we abandon this property and require only the existence of derivatives of the $k$-th order, we can extend regular and Dirac distributions on previously defined space $X^k$. Indeed,
\begin{align}
    | \dual{\partial_x^n \delta(x-y)}{f(x)} | \leq \norm{f}_n \quad \forall n = 0,1,\cdots, k\:,
\end{align}
so Dirac distributions and their derivatives are also bounded linear functionals on $X^k$ up to the order $k$.

Let $a \in C^\infty(\Omega)$. Then it makes sense to define the product of $a$ and arbitrary distribution $T$ by
\begin{align}
    \dual{a T}{\phi} = \dual{T}{a \phi} \:.
\end{align}

In the next chapter we shall frequently use the following identities of the delta distributions, which are written in the next lemma and which we will refer to as the Dirac identities.

\begin{lemma}[Dirac identities] \label{lemma:delta}
    \begin{align}
        f(y) \delta (x-y) =& f(x) \delta(x-y) \:, \\
        f(y) \partial_x \delta(x-y) =& f(x) \partial_x \delta(x-y) + f'(x) \delta(x-y) \:.
    \end{align}
\end{lemma}
\begin{proof}
    The first relation is trivial.
    Direct calculation gives
    \begin{align*}
        \dual{f(x) \partial_x \delta(x-y)}{\psi(x)} 
        =& \dual{\partial_x \delta(x-y)}{f(x) \psi(x)} 
         \\ =& -\dual{\delta(x-y)}{f'(x)\psi(x)} - \dual{\delta(x-y)}{f(x)\psi'(x)} 
         \\ =& -f'(y) \psi(y) - f(y) \dual{\delta(x-y)}{\psi'(x)} 
         \\ =& - \dual{f'(x)  \delta(x-y)}{\psi(x)} + f(y) \dual{\partial_x \delta(x-y)}{\psi(x)} 
         \\ =& \dual{f(y) \delta(x-y) - f'(x) \partial_x \delta(x-y)}{\psi(x)} \:.
    \end{align*}
    Therefore
    \begin{align}
        \dual{f(y) \partial_x \delta(x-y)}{\psi(x)} = \dual{f'(x) \delta(x-y)}{\psi(x)} + \dual{f(x) \partial_x \delta(x-y)}{\psi(x)} \:.
    \end{align} 
\end{proof}

The convolution of two distributions is not simple to construct in rigorous way. The procedure is the following. First the tensor product of distributions is defined
\begin{align}
    \dual{A(x) \otimes B(y)}{\phi(x,y)} = \dual{A(x)}{\dual{B(y)}{\phi(x,y)}} \:.
\end{align}
The product defined this way satisfies comutativity, associativity and it is indeed a distribution (altough the proof of these properties is non-trivial). Then the convolution of distribution is defined by
\begin{align}
    \dual{A \star B}{\phi(x)} = \dual{A(x) \otimes B(y)}{\phi(x+y)} \:,
\end{align}
with additional requirements on $A \otimes B$ and $\phi$. Convolution defined in this way satisfies associativity and if it exists, then
\begin{align}
    D^\alpha (A \star B) = D^\alpha A \star B = A \star D^\alpha B \:.
\end{align}

The convolution of Dirac distribution and arbitrary distribution $A$ always exists and is equal to
\begin{align}
    A(y) \star \delta(x-y) = A(x) \:.
\end{align}
This means, in particular, that given an integral functional $T[u]$ one can write
\begin{align}
    T[u(x)] = T[u(y)]  \star \delta(x-y)  \:,
\end{align}
or in physicist's notation
\begin{align}
    T[u(x)] = \int \D y \, T[u(y)] \, \delta(x-y) \:. \label{eq:konvoluce-s-deltou}
\end{align}
The equation \eqref{eq:konvoluce-s-deltou} is very important in calculations of functional derivatives.
\begin{example}
    Let $u(x)$ be a field variable and $F[u(x)]$ an integral functional. 
    To compute its functional derivative with respect to $u(y)$ in different position $y$, we use the integral expression of $F$
    \begin{align}
        F[u(x)] = \int \D y F[u(y)] \delta (x-y) \:.
    \end{align}
    Its Gateaux derivative is
    \begin{align}
        \delta F[u,h] = \der{}{t} F[u+th]|_{t=0} = \int \D y \pder{F}{u}(y) \delta(x-y) h(y) \:,
    \end{align}
    therefore
    \begin{align}
        \fder{F(u(x))}{u(y)} = \pder{F}{u}(x) \delta(x-y) \:.
    \end{align}
\end{example}

\begin{example}
    In the special case of performing functional derivative w.r.t. the same field but in the different position,
    \begin{align}
        \fder{u^i(x)}{u^j(y)} = \delta^i_j \delta(x-y) \:,
    \end{align}
    because
    \begin{align}
        \fder{u^i(x)}{u^j(y)} = \fder{}{u^j(y)} \int \D y u^i(y) \delta (x-y) = \pder{u^i}{u^j}(y) \delta (x-y) = \delta^i_j \delta (x-y) \:.
    \end{align}
    The same applies to
    \begin{align}
        \fder{}{u^j(y)} \partial_x u^i(x) = - \delta^i_j \partial_x \delta(x-y) \:,
    \end{align}
    because
    \begin{align}
        \fder{}{u^j(y)} \partial_x u^i(x) = \fder{}{u^j(y)} \int \D y \partial_y u^i(y) \delta(x-y) 
        \\ = -\fder{}{u^j(y)} \int \D y u^i(y) \partial_y \delta(x-y) = - \delta^i_j \partial_x \delta(x-y) \:.
    \end{align}
\end{example}
\chapter{Geometrical preliminarities}

In this chapter $M$ represents a finite-dimensional smooth manifold, $\mathcal F M$ the space of smooth functions $M \to \R$, $T_x M$ the tangent space in $x \in M$, $T M$ the tangent bundle on $M$, $T^p_q M$ the bundle of tensors of rank $(p,q)$ on $M$. We also automatically assume the so called Einstein summation over contravariant and covariant indices. Vectors, covectors and tensors of arbitrary rank are denoted (only in this chapter) by bold symbols, for the purpose of higliting the difference between the abstract indices (denoted by normal letters $i,j,\cdots$) and vector indices (denoted by bold letters $\vc A$, $\vc b$, $\cdots$)

\section{Covariant derivative. Torsion and curvature}

\begin{definition}
    A metric tensor $\vc g_{ij} \in \T^0_2 M$ is symmetric and non-degenerate, i.e.
    \begin{align}
        \vc g_{ij} =&\, \vc g_{ji} \quad \textbf{(symmetry)}\:, \\ \det \vc g_{ij} \neq&\, 0 \quad \textbf{(non-degeneracy)}\:.
    \end{align}
\end{definition}

If the metric $\vc g_{ij}$ is positive quadratic form, the manifold $M$ is said to be Riemannian. If it is indefinite, the manifold is said to be pseudo-Riemannian.

Non-degeneracy of the metric $\vc g_{ij}$ allows us to define the inverse metric $\vc g^{ij}$ defined by the property
\begin{align}
    \vc g^{ij} \vc g_{jk} = \vc \delta^i_k \:.
\end{align}

\begin{definition}[Covariant derivative]
    Let $\vc a \in TM$. A covariant derivative $\nabla_{\vc a}$ in the direction $\vc a$ is an operation on $T^p_q M$ satisfying these properties:
    \begin{align}
        &\nabla_{\vc a} (\vc A + r \vc B) = \nabla_{\vc a} \vc A + r \nabla_{\vc a} \vc B \quad \forall \vc A,\vc B \in T^k_l M \:,\: r \in \R \:, \quad \textbf{(linearity)} \\
        &\nabla_{\vc a} (\vc A \otimes \vc B) = (\nabla_{\vc a} \vc A) \otimes \vc B + \vc A \otimes (\nabla_{\vc a} \vc B) \quad \forall \vc A,\vc B \in T^k_l M \:, \quad \textbf{(Leibniz rule)} \\
        &\nabla_{\vc a} (\mathsf{C} \vc A) = \mathsf{C} (\nabla_{\vc a} \vc A) \quad \forall \vc A \in T^k_l M \:, \quad \textbf{(commutation with contraction)} \\
        &\nabla_{\vc a} f = \vc a[f] = \vc a \cdot \D f \quad \forall f \in \mathcal F M \:, \quad \textbf{(gradient on scalars)}
    \end{align}
    where $\mathsf{C}$ is the contraction operation.
\end{definition}
(Note that the covariant derivative $\nabla_{\vc a}$ has the lower index bold, as it refers to the vector $\vc a$.)

If we multiply the vector field $\vc a \in T M$ by some smooth function $f \in \mathcal F M$, then
\begin{align}
    \nabla_{f \vc a} \vc A = f \nabla_{\vc a} \vc A \quad \forall \vc A \in T^p_q M \quad \textbf{(ultralocality)}\:.
\end{align}
Therefore we can introduce the covariant differential $\nabla_j$ (with the normal lower index) as the operation from $\nabla : \T^p_q \to \T^p_{q+1}$ defined by the property
\begin{align}
    \nabla_{\vc a} \vc A = \vc a \cdot \nabla \vc A \:,
\end{align} 
or equivalently in the index notation
\begin{align}
    \nabla_{\vc a} \vc A^{ij\cdots}_{kl\cdots} = \vc a^k \nabla_k \vc A^{ij\cdots}_{kl\cdots} \:.
\end{align}

Any covariant diferential $\nabla_i$ can be expressed via the partial derivative with respect to the coordinates $(x^j)$ as
\begin{align}
    \nabla_i \vc A^{ab \dots}_{cd \dots} = \vc A^{ab \dots}_{cd \dots, i} + \vc \Gamma^{a}_{im} \vc A^{mb \dots}_{cd \dots} + \vc \Gamma^{b}_{im} \vc A^{am \dots}_{cd \dots} + \cdots
    - \vc \Gamma^{m}_{ic} \vc A^{ab \dots}_{md \dots} - \vc \Gamma^{m}_{id} \vc A^{ab \dots}_{cm \dots} - \cdots \:,
\end{align}
where $\vc \Gamma^{a}_{bc}$ are the affine connection coefficients (Christoffel symbols). Given two coordinate systems $(x^i)$ and $(\tilde x^a)$, the Christoffel symbols transform as
\begin{align}
    \vc {\tilde\Gamma}^{a}_{bc} = \vc {\Gamma}^i_{jk} \pder{x^j}{\tilde  x^b} \pder{x^k}{\tilde x^c}  \pder{\tilde x^a}{x^i} + \pder{\tilde x^a}{x^m} \ppder{x^m}{\tilde x^b}{\tilde x^c} \:.
\end{align}

\begin{definition}[Torsion and Riemann tensors]
    Let $[\cdot,\cdot]$ be the Lie bracket of two vector fields.

    \begin{enumerate}
        \item The torsion tensor $\vc T^{k}_{mn}$ is defined by
        \begin{align}
            \vc T^k_{mn} \vc a^m \vc b^n = \vc a^m \nabla_m \vc b^k - \vc b^m \nabla_m \vc a^k - [\vc a, \vc b]^k \quad \forall \vc a, \vc b \in TM \:.
        \end{align}
        \item The Riemann curvature tensor $\vc R\indices{_i_j^k_l}$ is defined by
        \begin{align}
            \vc R\indices{_i_j^k_l} \vc a^i \vc b^j \vc c^l = [\nabla_{\vc a} \nabla_{\vc b} - \nabla_{\vc b} \nabla_{\vc a} - \nabla_{[\vc a, \vc b]}] \vc c^k \quad \forall \vc a, \vc b, \vc c \in TM \:,
        \end{align}
    \end{enumerate}
\end{definition}


The torsion tensor represents the commutator of the second covariant diferentials applied on functions
\begin{align}
    [\nabla_i \nabla_j - \nabla_j \nabla_i] f = - \vc T^{m}_{ij} \D_m f \quad \forall f \in \mathcal F M \:.
\end{align}
and Riemann tensor represents the commutator on vectors
\begin{align}
    [\nabla_i \nabla_j - \nabla_j \nabla_i] \vc a^k =  - \vc T_{ij}^m \nabla_m \vc a^k + \vc R\indices{_i_j^k_m} \vc a^m \:. 
\end{align}

\begin{proposition}[Zero torsion and symmetry of the Christoffel symbols]
    The torsion tensor $\vc T^i_{jk}$ of the covariant derivative $\nabla_i$ is zero if and only if the corresponding Christoffel symbols are symmetric in the lower indices, i.e. $\vc \Gamma^i_{jk} = \vc \Gamma^i_{kj}$.
\end{proposition}


\begin{definition}[Levi-Civita covariant derivative]
    Let $M$ be a pseudo-Riemannian manifold with the metric $\vc g_{ij}$. Levi-Civita covariant derivative $\vc \nabla_k$ is defined by the following properties:
    \begin{align}
        \vc \nabla_k \vc g_{ij} = \vc g_{ij,k} - \vc \Gamma^m_{ki} \vc g_{mj} - \vc \Gamma^m_{kj} \vc g_{im}  =& 0 \quad \textbf{(compatibility with metric)} \:, \\
        \vc T^{i}_{jk} =& 0 \quad \textbf{(zero torsion)} \:,
    \end{align}
    where $\vc T^i_{jk}$ is the torsion corresponding to $\vc \nabla_j$.
\end{definition}

If $\vc g_{ij}$ is non-degenerate, i.e. $\det \vc g_{ij} \neq 0$, then the Levi-Civita covariant derivative is always defined and its Christoffel symbols are given by
\begin{align}
    \vc \Gamma^m_{ij} =\frac{1}{2} \vc g^{mn} \left( \vc g_{jn,i} + \vc g_{ni,j} - \vc g_{ij,n} \right) \:.
\end{align}

The covariant derivative compatible with the metric commutes with the operation of lowering and uppering indices of a tensor. For example:
\begin{align}
    \nabla_i \vc T^{ab} = \nabla_i (\vc g^{am} \vc T^b_m) = \vc g^{am} \nabla_i \vc T^b_m \:.
\end{align}

The Riemann curvature tensor for the Levi-Civita covariant derivative is
\begin{align}
    \vc R\indices{_i_j^k_l} = \vc \Gamma^k_{jl,i} - \vc \Gamma^k_{il,j} + \vc \Gamma^k_{im} \vc \Gamma^m_{jl} - \vc \Gamma^k_{jm} \vc \Gamma^m_{il} \:.
\end{align}

\section{Poisson brackets}

Poisson brackets and symplectic geometry play the essential role in the Hamiltonian mechanics and can be seen from many different perspectives. An inspirable book on this topic is e.g. \cite{Marsden}.  

\begin{definition}[Poisson bracket]
    A Poisson bracket is a bilinear operation $\set{\cdot, \cdot} : \F M \times \F M \to \F M$ which satisfies
    \begin{align}
        &\set{f,g} = -\set{g,f} \quad \textbf{(skew-symmetry)} \:, \\
        &\set{fg,h} = f \set{g,h} + g \set{f,h} \quad \textbf{(Leibniz identity)} \:, \\
        &\set{\set{f,g},h} + \set{\set{g,h},f} + \set{\set{h,f},g} = 0 \quad \textbf{(Jacobi identity)} \:.
    \end{align}
\end{definition}

\begin{definition}[Symplectic 2-form]
    A skew-symmetric 2-form $\vc \omega \in \T^0_2 M$ is said to be symplectic if it is closed and nondegenerate, i.e.
    \begin{align}
        &\D \vc \omega = 0 \quad \textbf{(closedness)} \:, \\
        &\det \vc \omega \neq 0 \quad \textbf{(nondegeneracy)}\:.
    \end{align}
    A pair $(M,\vc \omega)$ is called sympletic manifold.
\end{definition}

Since in odd dimension every skew-symmetric 2-form is degenerate, the dimension $m$ of the symplectic manifold $M$ must be even number. The following statement holds.

\begin{theorem}[Darboux theorem - symplectic forms formulation]
    Let $(M,\vc \omega)$ be a $2n$-dimensional symplectic manifold and $x \in M$. Then there is a neighbourhood $U(x)$ and a coordinate system $(p_1,\cdots,p_n,q^1,\cdots,q^n)$ such that $\vc \omega$ has the form
    \begin{align}
        \vc \omega = \D p_a \wedge \D q^a \quad \text{on } U(x) \:. \label{eq:standard form}
    \end{align}
\end{theorem}
We refer to this coordinate system as to the standard coordinates \\$(x^1,\cdots, x^{2n}) = (p_1,\cdots,p_n,q^1,\cdots,q^n)$ and to the form \eqref{eq:standard form} as to the standard form of $\vc \omega$.

\begin{definition}[Poisson bivector]
    Let $(M,\vc \omega)$ be a symplectic manifold.
    The Poisson bivector $\vc P \in \T^2_0 M$ is the inverse of $\vc \omega$, i.e.
    \begin{align}
        \vc P^{ij} \vc \omega_{jk} = \vc \delta^i_k \:.
    \end{align}
\end{definition}
Thus a symplectic manifold $(M,\vc \omega)$ has its associated Poisson bracket given by
\begin{align}
    \set{f,g} \coloneqq \vc P (\D f, \D g) = \vc P^{ab} \D_a f \D_b g \quad \forall f,g \in \F M \:. 
\end{align}
In the standard coordinates
\begin{align}
    \set{f,g} = \sum_{i=1}^n \left( \pder{f}{q^i} \pder{g}{p_i} - \pder{f}{p_i} \pder{g}{q^i} \right) \:.
\end{align}

In the standard coordinates we have
\begin{align}
    \set{x^i,x^j} =& \vc P^{ij} \:, \\
    \set{\set{x^i,x^j},x^k} =& \pder{\vc P^{ij}}{x^m} \pder{x^k}{x^n} \vc P^{mn} = \pder{\vc P^{ij}}{x^m} \vc P^{m k} \:,
\end{align}
so the Jacobi identity is given by
\begin{align}
    \pder{\vc P^{ij}}{x^m} \vc P^{m k} + \pder{\vc P^{jk}}{x^m} \vc P^{m i} + \pder{\vc P^{ki}}{x^m} \vc P^{m j} = 0 \:.
\end{align}

\begin{definition}[The Casimir of the Poisson bracket]
    The function $g \in \F M$ is called a Casimir for the given Poisson bracket $\set{\cdot,\cdot}$, if
    \begin{align}
        \set{f,g} = 0 \quad \forall g \in \F M \:.
    \end{align}
\end{definition}

\chapter{The connection between Riemannian geometry and Hamiltonian mechanics}

In this chapter we shall prove the theorem of the connection between the hydrodynamic Poisson bracket and the metrics found in the homogeneous system of hydrodynamic type. The original result was given by Dubrovin and Novikov and can be found in \cite{Dubrovin-Novikov}. The proof itself and the preceding lemmas are based on a considerable amount of technical computations, which we shall clarify in detail in order to offer the reader a comprehensive understanding of the chapter in question.





\section{Systems of hydrodynamic type}

\begin{definition}[Homogeneous system of hydrodynamic type]
    A homogeneous system of hydrodynamic type is an equation of the form
    \begin{align}
        \label{eq:h-system}
        \pder{u^i}{t} = f_j^{i \alpha} (u) \pder{u^j}{x^\alpha} 
    \end{align}
    where $u^i(t,x^\alpha)$ are unknown functions, $i = 1, \dots , N$, $\alpha = 1, \dots , d$.
\end{definition}

Sometimes we refer to the number of equations $N$ as the dimension of the hydrodynamical phase space and to the number of spatial variables $d$ as the dimension of the configuration space.

Before we go any further, let us present a few examples.

\begin{example}
    The energy is given by the functional
    \begin{align}
        E = \int e(\rho, \vc m, s, \vc F) \, \D \vc x
    \end{align}
    The evolution equations in the Eulerian frame are
    \begin{align}
        \partial_t \rho =& - \partial_i (\rho E_{m_i}) 
        \\
        \partial_t m_i =& - \partial_j(m_i E_{m_j}) - \rho \partial_i E_\rho - m_j \partial_i E_{m_j} - m_j \partial_i E_{m_j} 
        \\
        -& s \partial_i E_s - F^j_J \partial_i E_{F^j_J} + \partial_j(F^j_I E_{F^i_I} + F^i_O E_{F^i_I}) 
        \\
        \partial_t s =& - \partial_i (sE_{m_i})
        \\
        \partial_t F^i_I =& - E_{m_k} \partial_k F^i_I + F^j_I \partial_j E_{m_i} \:.
    \end{align}
    Total energy density $e$ can be prescribed as
    \begin{align}
        e = \frac{\vc m^2}{2 \rho} + \epsilon(\rho,s , \vc F) \:,
    \end{align}
    where $\epsilon$ is the elastic and internal energy. In particular, $E_{\vc m} = \frac{\vc m}{\rho} = \vc v$ becomes the velocity. The evolution equation for the deformation gradient then gets the explicit form
    \begin{align}
        \partial_t \vc F = -(\vc v \cdot \nabla) \vc F + \nabla \vc V \cdot \vc F \:,
    \end{align}
    which is the usual evolution equation for $\vc F$ in the Eulerian frame.
\end{example}

It was Riemann who noticed that the functions $f^{i \alpha}_j$ in \eqref{eq:h-system} are in fact tensors [CITACE]\cite{Dubrovin-Novikov}.

\begin{proposition}[Transformation of $f^{i \alpha}_j$] 
    \label{prop:transformace-A} 
    Under smooth change of variables $u^i \mapsto v^a$ of the form
    \begin{align}
        u^i = u^i(v^1,\cdots, v^N)
    \end{align}
    the functions $f^{i\alpha}_{j}$ transform for each $\alpha$ according to the tensor law
    \begin{align}
        f^{a \alpha}_{b}(v) = \pder{v^a}{u^i} \pder{u^j}{v^b} f^{i \alpha}_j (u) \:.
    \end{align}
\end{proposition}
\begin{proof}
    Using a direct calculation we obtain
    \begin{align}
        \pder{u^i}{t}(v) = \pder{u^i}{v^a} \pder{v^a}{t} \:, \\
        \pder{u^j}{x^\alpha}(v) = \pder{u^j}{v^b} \pder{v^b}{x^\alpha}\:,
    \end{align}
    so
    \begin{align}
        \pder{u^i}{v^a} \pder{v^a}{t} = \pder{u^j}{v^b} f^{i \alpha}_j (v(u)) \pder{v^b}{x^\alpha}
    \end{align}
    and therefore
    \begin{align}
        \pder{v^a}{t} = \underbrace{\pder{v^a}{u^i} \pder{u^j}{v^b} f^{i \alpha}_j(v(u))}_{f^{a \alpha}_b(v)} \pder{v^b}{x^\alpha} \:.
    \end{align}
\end{proof}

According to the previous lemma, if we denote the space of points $(t,x^\alpha) \subset \R \times \R^N$ as some manifold $M$, then the unknown variables $u^i(t,x^\alpha)$ can be seen as the local coordinates of that manifold. The functions $f^{i \alpha}_j$ are transformed as tensors of the type (1,1) in the $i$ and $j$ indices.

Our next purpose is to gain a richer geometry of that space $M$. If the functions $f^{i \alpha}_j$ have a special form, we can introduce a geometry of the Poisson brackets.

\begin{definition}
    \begin{enumerate}
        \item A Poisson bracket of the hydrodynamic type of two functionals $I_1$ a $I_2$ is defined as
        \begin{align}
            \set{I_1,I_2} = \int \D x 
            \left[ \frac{\delta I_1}{\delta u^q(x)} A^{qp} \frac{\delta I_2}{\delta u^p(x)} \right]
        \end{align}
        where
        \begin{align}
            A = (A^{qp}) = \left( g^{qp \alpha}(u) \der{}{x^\alpha} + b^{qp\alpha}_s(u) \pder{u^s}{x^\alpha} \right) \:.
        \end{align}
        where $g^{ij \alpha}$ and $b_k^{ij \alpha}$ are certain functions, $i,j,k = 1, \dots, N$ and $\alpha = 1, \dots, d $.

        In particular
        \begin{align}
            \set{u^i(x), u^j(y)} = g^{ij \alpha} [u(x)] \delta_\alpha(x-y) + b_k^{ij \alpha} [u(x)] u_\alpha^k(x) \delta(x-y) \:.
        \end{align}
        
        \item The functional of hydrodynamic type is a functional of the form
        \begin{align}
            H[u] = \int_{\R^d} h(u(x)) \D^d x \:,
        \end{align}
        where $h$ is independent of $u_\alpha$, $u_{\alpha \beta}$. The function $h$ is called a hamiltonian density.

        \item The hamiltonian system of hydrodynamic type is a system of the form
        \begin{align}
            u^i_t (x) = 
            \left[ g^{ij \alpha} [u(x)] \ppder{h(u)}{u^j}{u^k} + b_k^{ij \alpha} [u(x)] \pder{h(u)}{u^j} \right] u_\alpha^k(x) \:,
        \end{align}
        where $H[u]$ is the functional of the hydrodynamic type and $h$ the hamiltonian density.
    \end{enumerate}
\end{definition}

\begin{proposition}
    \begin{align}
        u^i_t(x) = \set{u^i,H} \:.
    \end{align}
\end{proposition}
\begin{proof}
    \begin{align}
        \set{u^i(x),H} = \int \D x' \D y' \fder{u^i(x)}{u^p(x')} \left[ g^{pq}(u(x')) \der{}{x'} + b^{pq}_m \pder{u^m}{x} \right] \fder{H}{u^q(y')} \:.
    \end{align}
    \begin{align}
        \fder{u^i(x)}{u^p(x')} = \delta^i_p \delta(x-x')
    \end{align}
    \begin{align}
        H[u] = \int h(u(y')) \D y' \:, \quad \fder{H}{u^q(y')} = \int \pder{h}{u^q}(y) \delta(y-y')
    \end{align}
    \begin{align}
        \set{u^i(x),H} =& \int \D y' \left[ g^{iq}(u(x)) \der{}{x} + b^{iq}_m(u(x)) \pder{u^m}{x} \right] \pder{h}{u^q}(y) \delta(y-y') 
        \\ =& g^{iq}(u(x)) \pder{^2 h}{u^k \partial u^q} \pder{u^k}{x} + b^{iq}_m(u(x)) \pder{u^m}{x} \pder{h}{u^q}
        \\ =& \left[ g^{ij}(u(x)) \pder{^2 h}{u^j \partial u^k}  + b^{ij}_k(u(x))  \pder{h}{u^j} \right] \pder{u^k}{x} \:.
    \end{align}
\end{proof}
\todo[inline]{Dokončit a přidělat $\alpha$. kde se předělá $\der{}{x}$ na $\der{}{y}$?}

Note that we have yet to show that the Poisson bracket of the hydrodynamic type defined above is indeed a Poisson bracket in the classical sense. We must verify that it satisfies the three properties: skew-symmetry, Leibniz rule and Jacobi identity. In fact, this will be the key theorem of this chapter.

\section{One-dimensional case}

In this section we first consider the one-dimensional case in which the unknown functions $u^i$ depend on only two variables $(t,x)$. The hamiltonian system is then
\begin{align}
    \pder{u^i}{t} (t,x) =&
    \left[ g^{ij} [u(t,x)] \ppder{h(u)}{u^j}{u^k} + b_k^{ij} [u(t,x)] \pder{h(u)}{u^j} \right] \pder{u^k}{x}(t,x) 
\end{align}
and the Poisson bracket of hydrodynamic type is of the form
\begin{align}
    \set{I_1,I_2} = \int \D x
    \left[ \frac{\delta I_1}{\delta u^q(x)} A^{qp} \frac{\delta I_2}{\delta u^p(x)} \right] \:,\quad \text{where } A^{qp} = \left( g^{qp}(u) \der{}{x} + b^{qp}_s(u) \pder{u^s}{x} \right) \:.
\end{align}
In particular
\begin{align}
    \label{eq:Poisson1D}
    \set{u^i(x), u^j(y)} =& g^{ij} [u(t,x)] \partial_x \delta(x-y) + b_k^{ij} [u(t,x)] \pder{u^k}{x} \delta(x-y) \:. 
\end{align}

We say that the Poisson bracket is non-degenerate if $\det(g^{ij}) \neq 0$. We will show later that the non-degeneracy is invariant under the transformation of the coordinates.

\begin{definition}
    Let $(g^{ij}(u))$ be a non-degenerate matrix. We define functions $\Gamma^{i}_{jk}(u)$ by
        \begin{align}
            b_k^{ij}(u) = -g^{is}(u) \Gamma^j_{sk}(u) \:, \quad i,j,k,s=1,\dots,N \:. \label{eq:definice-Gamma}
        \end{align}
\end{definition}

The choice of denotation $g^{ij}$ and $\Gamma^{i}_{jk}$, respectively, has its purpose. We will show soon that these are in fact the metric and afinne coeficients, respectively, on $M$.

\begin{proposition}[Leibniz rule for Poisson bracket]
    The Poisson brackets of the hydrodynamic type satisfy the Leibniz rule
    \begin{align}
        \label{eq:transformace-zavorka}
        \set{v^p(u^i(x)),v^q(u^j(y))} = \pder{v^p}{u^i}(x) \pder{v^q}{u^j}(y) \set{u^i(x),u^j(y)} \:. 
    \end{align}
\end{proposition}
\begin{proof}
    The coordinates are written as functionals in the form
    \begin{align}
        v^p(u^i(x)) = \int \D x' v^p(u^i(x')) \delta (x-x') \:.
    \end{align}
    Applying this to the definition of the Poisson bracket \label{eq:Poisson1D}, we have
    \begin{align}
        \set{v^p(u^i(x)),v^q(u^j(y))} 
        =& \int \D x' \int \D y' \underbrace{\fder{v^p}{u^i(x')} }_{\pder{v^p}{u^i}(x') \delta(x-x')} \set{u^i(x'),u^j(y')} \underbrace{\fder{v^q}{u^j(y')} }_{\pder{v^q}{u^j}(y') \delta(y-y')}
        = \\ =& \pder{v^p}{u^i}(x) \pder{v^q}{u^j}(y) \set{u^i(x),u^j(y)} \:.
    \end{align}
\end{proof}


\begin{proposition}[Transformation of $g$ and $\Gamma$] \label{prop:transoformace}
    Let $u^i \mapsto v^a$ be the smooth change of variables. Then
    \begin{enumerate}
        \item The coeficients $g^{ij}$ transform as the (2,0) tensors, i.e. 
        \begin{align}
            g^{pq}(v) = \pder{v^p}{u^i} \pder{v^q}{u^j} g^{ij}[u(v)] \:, \quad p, q = 1, \dots, N \:. \label{eq:transformace-metrika}
        \end{align}
        \item The coeficients $\Gamma^i_{jk}$ transform like the afinne connection symbols, i.e.
        \begin{align}
            \Gamma^p_{qr} (v) = \pder{v^p}{u^i} \pder{u^j}{v^q} \pder{u^k}{v^r} \Gamma^i_{jk}(u) + \pder{v^p}{u^i} \ppder{u^i}{v^q}{v^r} \:. \label{eq:transformace-konexe}
        \end{align}
    \end{enumerate}
\end{proposition}

\begin{proof}
    \begin{enumerate}
        \item Let us apply the transformation equation \eqref{eq:transformace-zavorka} to the definition of the Poisson bracket \eqref{eq:Poisson1D}. We get 
        \begin{align}
            g^{pq}[v(u(x))] \partial_x \delta(x-y) + b^{pq}_s[v(u(x))] \pder{v^s}{u^k}(x) \pder{u^k}{x} \delta(x-y) \nonumber
            \\ = \pder{v^p}{u^i}(x) \pder{v^q}{u^j}(y) \left[ g^{ij}[u(x)] \partial_x \delta(x-y) + b^{ij}_k[u(x)] \pder{u^k}{x} \delta(x-y) \right] \:. \label{eq:dosad}
        \end{align}
        From now we denote for readibility
        \begin{align}
            T^a_k(x) := \pder{v^a}{u^k}(x) \:.
        \end{align}
        On the right-hand side of \eqref{eq:dosad} we apply Dirac identity (Lemma \vref{lemma:delta}) on $T^q_j(y) \partial_x \delta(x-y)$. All functions are then of variable $x$, so from now we won't write it explicitely. We have
        \begin{align}
            g^{pq}(v) \partial_x \delta(x-y) + b^{pq}_n(v) T^n_j \pder{u^j}{x} \delta(x-y) \nonumber
            \\ = T^p_i T^q_j g^{ij}(u) \partial_x \delta(x-y) 
            + \left( \ppder{v^q}{u^j}{u^k} + T^p_i T^q_j b^{ij}_k \right) \pder{u^k}{x} \delta(x-y) \:. \label{eq:porovnani}
        \end{align}
        The terms with $\partial_x \delta$ must be equal to each other, as must $\delta$ terms. From $\partial_x \delta$ terms we get
        \begin{align}
            g^{pq}(v) = T^p_i T^q_j g^{ij}(u) \:,
        \end{align}
        which is the transformation relation for $g$.

        \item From $\delta$ terms in \eqref{eq:porovnani}, we get
        \begin{align}
            b^{pq}_n(v) T^n_k =
            b^{ij}_k(u) T^p_i T^q_j + g^{ij}(u) T^p_i \ppder{v^q}{u^i}{u^k}
        \end{align}
        and after \eqref{eq:definice-Gamma}
        \begin{align}
            - g^{ps}(v) \Gamma^{q}_{sn}(v) T^n_k =
            - g^{il}(u) \Gamma^j_{lk}(u) T^p_i T^q_j + g^{ij}(u) T^p_i \ppder{v^q}{u^i}{u^k} \:.
        \end{align}
        On the left-hand side we use the transformation of $g$ \eqref{eq:transformace-metrika}
        \begin{align}
            -g^{ij}(u) \Gamma^q_{sm}(v) T^p_i T^s_j T^m_k = - g^{ij}(u) \Gamma^l_{jk} T^p_i T^q_l + g^{ij}(u) \ppder{v^q}{u^i}{u^k} T^p_i 
        \end{align}
        and after vectoring out $g^{ij}(u) T^p_i$
        \begin{align}
            \Gamma^l_{jk}(u) T^q_l = \Gamma^q_{sm}(v) T^s_j T^m_k + \ppder{v^q}{u^i}{u^k} \:.
        \end{align}
        All that remains is to multiply this equation with $(T^{-1})^a_q$ and we get the final transformation relation for $\Gamma$
        \begin{align}
            \Gamma^a_{jk}(u) = \Gamma^q_{sm}(v) \pder{v^s}{u^j} \pder{v^m}{u^k} \pder{u^a}{v^q} + \ppder{v^q}{u^i}{u^k} \pder{u^a}{v^q} \:.
        \end{align}
    \end{enumerate}
\end{proof}

Let us summarize the results. We have defined Poisson bracket in the way that it satisfies the Leibniz rule. It remains to investigate the skew-symmetry and Jacobi identity. We have also defined $g^{pq}$ and $\Gamma^p_{qr}$, which transform like tensors and Christoffel symbols, respectively. It remains to show that $g^{pq}$ is indeed symmetric and $\Gamma$ is the Levi-Civita connection, i.e. the corresponding covariant derivative annuls the metric. We are now ready to prove the key theorem of this chapter which shows that these properties are indeed equivalent. Moreover, the corresponding covariant derivative has zero curvature, which indicates that the phase space $M$ is in fact Euclidean space.

\begin{theorem}[Dubrovin, Novikov] \label{theorem1:ekvivalence}
    Let $\det(g^{ij}) \neq 0$. Then \eqref{eq:Poisson1D} defines a Poisson bracket with skew-symmetry, Leibniz rule and Jacobi identity if and only if these three conditions are satisfied:
    \begin{enumerate}
        \item $g^{ij} = g^{ji}$, that is, $g^{ij}$ is a metric on the phase space $M$.
        \item $\Gamma^{i}_{jk} = \Gamma^i_{kj}$, that is, $\Gamma^{i}_{jk}$ are affine coeficients, such that the corresponding covariant derivative has zero torsion and anihilate the metric $g^{ij}$.
        \item The covariant derivative corresponding to $\Gamma^{i}_{jk}$ has zero curvature.
    \end{enumerate}
\end{theorem}

\begin{corollary}
    There are local coordinates $w^i = w^i(u^1, \dots, u^N)$, $i = 1, \dots, N$ such that
    \begin{align}
        g^{ij}(w) = \tilde g^{ij} = \const \:, \quad \Gamma^i_{jk} (w) = 0 \:.
    \end{align}
    In these coordinates is Poisson bracket reduced to the constant one
    \begin{align}
        \set{w^i(x),w^j(y)} = \tilde{g}^{ij} \partial_x \delta(x-y) \:.
    \end{align}
    The only local invariant of the Poisson bracket is the signature of the metric $g^{ij}(u)$.
    \todo[inline]{Co se tímhle myslí? Pořád nemám rozmyšleno.}
\end{corollary}

\begin{proof}[Proof of theorem \ref{theorem1:ekvivalence}]
    The proof is divided into three steps.
    \\
    \textbf{Step 1: Skew-symmetry of Poisson bracket implies the symmetry of $g^{ij}$ and compatibility with $\Gamma^i_{jk}$.}

    Let the Poisson bracket be skew-symmetric, i.e.
    \begin{align}
        \set{u^i(x),u^j(y)} + \set{u^j(y),u^i(x)} = 0 \:.
    \end{align}
    The second bracket reads as
    \begin{align}
        \set{u^j(y),u^i(x)} =  g^{ji}[u(y)] \partial_x \delta(y-x) + b^{ji}_k[u(y)] \pder{u^k}{y} \delta(y-x) \:.
    \end{align}
    Let us apply Dirac identites (Lemma \vref{lemma:delta}) on $g^{ji}[u(y)]\partial_x \delta(y-x)$. We get
    \begin{align}
        \set{u^j(y),u^i(x)} =& \nonumber \\ 
        = - g^{ji} [u(x)] \partial_x \delta(x-y) -& \pder{g^{ji}}{u^k}[u(x)] \pder{u^k}{x} \delta(x-y) + b^{ji}_k[u(x)] \pder{u^k}{x} \delta(x-y)
        \:.
    \end{align}
    The skew-symmetry is then
    \begin{align}
        0 = \set{u^i(x),u^j(y)} + \set{u^j(y),u^i(x)} = \nonumber
        \\= \left[ g^{ij}[u(x)]  - g^{ji}[u(x)] \right] \partial_x \delta(x-y) + \left[ b^{ij}_k - \pder{g^{ji}}{u^k}  + b^{ji}_k  \right] \pder{u^k}{x} \delta(x-y) \:.
    \end{align}
    The right-hand side is equal to zero if and only if
    \begin{align}
        g^{ij}(u) = g^{ji}(u) \:, \label{eq:Th.Novikov-symetrie} \\
        \pder{g^{ij}}{u^k}(u) = b^{ij}_k(u) + b^{ji}_k(u)  \:. \label{eq:Th.Novikov-konexe}
    \end{align}
    Equation \eqref{eq:Th.Novikov-symetrie} says that $g^{ij}$ is a symmetric tensor. By assumption $\det g^{ij}$ is non-degenerate, so it defines a metric structure on $M$. Equation \eqref{eq:Th.Novikov-konexe} is equivalent to
    \begin{align}
        \pder{g^{ij}}{u^k} + g^{is} \Gamma_{sk}^j + g^{sj} \Gamma_{sk}^i = 0 \:.
    \end{align}
    If $\nabla$ is the covariant derivative corresponding to $\Gamma^i_{jk}$, then $\nabla_k g^{ij} = 0$. So $\Gamma_{ij}^k$ are affine connection coeficients and are compatible with $g^{ij}$.

    \textbf{Step 2: Jacobi identity is equivalent to zero torsion and curvature.} 
    
    Let us assume that Jacobi identity holds. We denote
    \begin{align}
        \mathcal J^{ijk}(x,y,z) = J^{ijk}(x,y,z)+ J^{jki}(y,z,x)+ J^{kij}(z,x,y) = 0
    \end{align}
    where
    \begin{align}
        J^{ijk}(x,y,z) = \set{u^k(z) , \set{u^i(x),u^j(y)}} \:.
    \end{align}
    \textbf{[TADY TO JEŠTĚ NENÍ HOTOVÉ.]}
    From now on we shall use the abbreviated notation
    \begin{align}
        g^{ij}_{,m} = \pder{g^{ij}}{u^m} \:, \quad b^{ij}_{k,m} = \pder{b^{ij}_k}{u^m} \:, \quad u^k_x = \pder{u^k}{x}  \:.
    \end{align}

    \begin{align*}
        J^{ijk}(x,y,z) =
        \int \D \xi \fder{u^k(z)}{u^n(\xi)}
        \left[ g^{mn}(\xi) \partial_\xi + b^{mn}_p \partial_\xi u^p \cdot \right] \fder{\set{u^i(x),u^j(y)}}{u^m(\xi)}  \:.
    \end{align*}
    The functional derivatives are
    \begin{align}
        \fder{u^k(z)}{u^n(\xi)} =& \delta^k_n \delta(z-\xi) \:, \\
        \fder{\set{u^i(x),u^j(y)}}{u^m(\xi)} =& \fder{}{u^m(\xi)} \left[ g^{ij}(x) \partial_x \delta(x-y) + b^{ij}_k(x) \partial_x u^k \delta(x-y)\right] = 
        \nonumber
        \\ =& \left[ g^{ij}_{,m} \partial_x \delta(x-y) + b^{ij}_{,m}(x) u^k_x \delta(x-y) \right] \delta (x-\xi) \:.
    \end{align}
    Thus we can eliminate the integration in $\xi$ and get rid of the integral.
    \begin{align*}
        & J^{ijk}(x,y,z) = \int \D \xi \delta^k_n \delta(z-\xi) \left[ g^{mn}(\xi) \partial_\xi + b^{mn}_p \partial_\xi u^p \cdot \right] \\ 
        & \left[ g^{ij}_{,m}(x) \partial_x \delta(x-y) \delta(x-\xi) + b^{ij}_{k,m}(x) u^k_x \delta(x-y) \delta(x-\xi) \right] 
        \\
        =& \left[ g^{mk}(\xi) \partial_z + b^{mk}_p \partial_z u^p \cdot \right] \left[ g^{ij}_{,m}(x) \partial_x \delta(x-y) \delta(x-\xi) + b^{ij}_{k,m}(x) u^k_x \delta(x-y) \delta(x-\xi) \right]
        \\
        =& \left[ g^{ij}_{,m} g^{mk} \partial_x u^m \partial_x \delta(x-y) + b^{ij}_{k,m} g^{mk} u^k_x \delta(x-y) \right] \partial_x \delta(z-x) + 
        \\
        +& \left[ g^{ij}_{,m} b^{mk}_p u^p_x \partial_x \delta(x-y) + b^{ij}_{k,m} b^{mk}_p u_x^p  u_x^k \delta(x-y) \right] \delta(z-x) \:.
    \end{align*}

    %\begin{align*}
    %    & J^{ijk}(x,y,z) = \\
    %    & \left[ g^{ij}_{,m}(x) \partial_x \delta(x-y) + b^{ij}_{k,m}(x) u^k_x \delta(x-y) \right] \left[ g^{mk}(x) \partial_x \delta(z-x) + b^{mk}_p(x) u^p_x \delta(z-x)\right] \:.
    %    \\ 
    %    =& \left[ g^{ij}_{,m} g^{mk} \partial_x u^m \partial_x \delta(x-y) + b^{ij}_{k,m} g^{mk} u^k_x \delta(x-y) \right] \partial_x \delta(z-x) + 
    %    \\
    %    +& \left[ g^{ij}_{,m} b^{mk}_p u^p_x \partial_x \delta(x-y) + b^{ij}_{k,m} b^{mk}_p u_x^p  u_x^k \delta(x-y) \right] \delta(z-x) \:.
    %\end{align*}


    We get four terms which are $\delta$ distributions and their derivatives.
    \begin{align*}
        J^{ijk}(x,y,z) =
        g^{ij}_{,m} g^{mk} \partial_x \delta(x-y) \partial_x \delta(z-x) 
        + b^{ij}_{k,m} g^{mk} u^k_x \delta(x-y) \partial_x \delta(z-x) 
        \\+ g^{ij}_m b^{mk}_p u^p_x  \partial_x \delta(x-y) \delta(z-x) 
        + b^{ij}_{k,m} b^{mk}_p u^p_x u^k_x \delta(x-y) \delta(z-x) \:.
    \end{align*}
    Jednotlivé členy označíme \begin{align*}
        J^{ijk}(x,y,z) = E(x) \partial_x \delta(x-y) \partial_x \delta(z-x) + F(x) \delta (x-y) \partial_x \delta(z-x) 
        \\+ G(x) \partial_x \delta(x-y) \delta(z-x) + H(x) \delta(x-y) \delta(z-x) \:.
    \end{align*}
    Díky symetrii, resp. antisymetrii $\delta$ funkcí
    \begin{align}
        \delta (x-y) = \delta(y-x) \:, \quad \partial_x \delta(x-y) = - \partial_y \delta(y-x)
    \end{align}
    Lze členy předělat na tvar
    \begin{align}
        J^{ijk}(x,y,z) = E(x) \partial_y \delta(y-x) \partial_z \delta(x-z) - F(x) \delta(y-x) \partial_z \delta(x-z) \\
        - G(x) \partial_y \delta(y-x) \delta(x-z) + H(x) \delta(x-y)  \delta(x-z)
    \end{align}


    \textbf{[ODTEĎ UŽ JE TO DODĚLANÉ.]}
    One of these equations reads as
    \begin{align}
        A_{02}^{ijk} = b^{ij}_s g^{sk} - b^{kj}_s g^{si} = -g^{ip} \Gamma^{j}_{ps} g^{sk} + g^{km} \Gamma^{j}_{ms} g^{si} =0 \:,
    \end{align}
    therefore, after multiplying with $g_{ai}g_{bk}$
    \begin{align}
        \Gamma^j_{ab} - \Gamma^j_{ba} = 0 \:.
    \end{align}
    So the corresponding covariant derivative has zero torsion.

    Another equation reads as
    \begin{align}
        A_{00}^{ijk} = B^{ijk}_p(u) u^p_{xx} + C^{ijk}_{pq}(u) u^p_x u^q_x = 0 \:,
    \end{align}
    where
    \begin{align}
        B^{ijk}_p 
        = (b^{jk}_{s,p} - b^{jk}_{p,s}) g^{si} + b_s^{ij} b_p^{sk} - b_s^{ik} b_p^{sj} \:.
    \end{align}
    The function $B^{ijk}_p$ beeing zero implies the zero curvature, which follows from
    \begin{align}
        B^{ijk}_p  
        =& -g^{is} \left( g^{jm}_{,p} \Gamma^k_{ms} + g^{jm} \Gamma^k_{ms,p} - g^{jm}_{,s} \Gamma^k_{mp} - g^{jm} \Gamma^k_{mp,s} \right) 
        \\ +& g^{im} \Gamma^j_{ms} g^{sn} \Gamma^{k}_{np} - g^{im} \Gamma^k_{ms} g^{sn} \Gamma^j_{np}
        = \\ =& g^{is} g^{jn} \left( \Gamma^m_{np} \Gamma^k_{ms} - \Gamma^{m}_{ns} \Gamma^k_{mp} + \Gamma^{k}_{mp,s} -  \Gamma^{k}_{ms,p} \right) 
        \\ +& g^{is} g^{nm} \left( \Gamma^j_{np} \Gamma^{k}_{ms} - \Gamma^j_{ns} \Gamma^k_{mp} \right) + g^{im} g^{sn} \left( \Gamma^j_{ms} \Gamma^k_{np} - \Gamma^k_{ms} \Gamma^j_{np} \right)
        = \\ =& -g^{is} g^{jn} \left( \Gamma^{m}_{ns} \Gamma^k_{mp}  - \Gamma^m_{np} \Gamma^k_{ms} + \Gamma^{k}_{ms,p} - \Gamma^{k}_{mp,s} \right)
        = \\ =& -g^{is} g^{jn} R^k_{nps} \:.
    \end{align}
    So far we proved that if hydrodynamical Poisson bracket satisfies skew-symmetry and Jacobi identity, the conditions 1-3 are satisfied.

    \textbf{Step 3:}

    If $\nabla$ has zero torsion and curvature, there are coordinates $w^i=w^i(u^1,\dots,u^N)$ for $i=1,\dots,N$ such that $g^{ij} = \const$ a $b^{ij}_k = 0$. In these coordinates
    \begin{align}
        \set{w^i(x),w^j(y)} = \tilde{g}^{ij} \partial_x \delta(x-y) \:.
    \end{align}
    Jacobi identity, skew-symmetry and Leibniz rule are for that bracket satisfied trivially. Therefore the three conditions are also sufficient. That finishes the proof.
\end{proof}

\subsection{Kdy je hydrodynamický systém hamiltonovský?}

Nyní můžeme explicitně přepsat podmínky, při kterých je obecný systém hydrodynamických rovnic hamiltonovský vůči nějaké nedegenerované Poissonově závorce. Nejprve si povšimneme, že funkce $f^i_j(u)$ lze přepsat pomocí Laplaceova–Beltramiho operátoru.

\begin{proposition}[O zápisu pomocí Laplaceova–Beltramiho operátoru]
    Mějme hamiltonovský systém rovnic
    \begin{align}
        \pder{u^i}{t} = f^i_k(u) \pder{u^k}{x} \:, \quad f^i_k(u) = g^{ij}(u) \ppder{h}{u^j}{u^k} - g^{is}(u) \Gamma^j_{sk}(u) \pder{h}{u^j} \:. \label{eq:hamiltonovsky-system}
    \end{align}
    Pak lze psát 
    \begin{align}
        f^i_k(u) = \nabla^i \nabla_k h(u) \:, 
    \end{align}
    kde $\nabla_j$ je Levi-Civitova kovariantní derivace metriky $g_{ij}$ a $\nabla^i = g^{is} \nabla_s$. 
\end{proposition}
\begin{proof}
    Přímým dosazením
    \begin{align}
        \nabla_k h(u) =& \pder{h}{u^k} \:, \\
        \nabla^i \nabla_k h(u) =& g^{is} \nabla_s \pder{h}{u^k} = g^{is} \left( \ppder{h}{u^s}{u^k} - \Gamma^j_{sk} \pder{h}{u^j} \right) = f^i_k(u) \:.
    \end{align}
\end{proof}

Pomocí tohoto zápisu snadno můžeme sepsat postačující podmínky hamiltonovskosti.
Důkaz plyne ihned z teorému \vref{theorem1:ekvivalence}.

\begin{theorem}[Postačující podmínka hamitonovskosti systému]
    Systém $u^i_t = f^i_j(u) u^j_x$ je hamiltonovský právě tehdy, když existuje nedegenerovaná metrika $g^{ij}(u)$, jejíž afinní konexe má nulovou křivost a splňuje
    \begin{align}
        g_{ij} f^k_j = g_{jk} f^k_i \:, \label{eq:gf=gf} \\
        \nabla_i f^k_j = \nabla_j f^k_i \label{eq:nablaf} \:.
    \end{align}
\end{theorem}

Speciálně vztah \eqref{eq:nablaf} říká, že $\nabla_i$ má nulovou torzi.


\subsection{Rekonstrukce metriky z $f^i_j(u)$?}

Máme-li zadaný hamiltonovský systém s maticí $f^i_j(u)$, lze metriku $g^{ij}(u)$ zkonstruovat jednoznačně? Tuto otázku nyní vyřešíme pro $N \geq 3$. 

Označme $\lambda_\alpha$ vlastní čísla matice $f^i_j(u)$. (Mohou být komplexní.) Předpokládejme, že jsou navzájem různá. Označme odpovídající bázi vlastních vektorů $e_\alpha(u)$. Definujme koeficient $c^\gamma_{\alpha \beta}(u)$ vztahem
\begin{align}
    [e_\alpha,e_\beta] = c^\gamma_{\alpha \beta} e_\gamma \:,
\end{align}
kde $[\cdot,\cdot]$ značí obyčejný komutátor funkcí. Předpokládejme dále, že pro navzájem různé $\alpha, \beta, \gamma$ je $c^\gamma_{\alpha \beta}$ různé od nuly .

\begin{definition}
    Matici $f^i_j(u)$ splňující podmínky výše nazveme hamiltonovskou maticí.
\end{definition}

\begin{theorem}[O rekonstrukci metriky]
    Nechť $N \geq 3$. Nechť je dána hamiltonovská matice $f^i_j(u)$. Pak lze zkonstruovat nedegenerovanou metriku $g^{ij}(u)$ s nulovou křivostí jednoznačně až na násobek konstantou.
\end{theorem}

\todo[inline]{Tomuhle nerozumím a myslím si, že v článku mají typo.}
\begin{proof}
    Z rovnice \eqref{eq:gf=gf} je vidět, že v bázi $e_\alpha$ je metrika $g^{ij}$ diagonální. 
    \todo[inline]{Opravdu?}
    V této bázi pak bude mít rovnice \eqref{eq:nablaf} tvar (zde se nesčítá přes opakované indexy)
    \begin{align}
        \partial_\alpha \lambda_\beta \delta_\beta^\gamma - \partial_\beta \lambda_\alpha \delta_\alpha^\gamma + (\Gamma_{\alpha \beta}^\gamma - \Gamma_{\beta \alpha}^\gamma ) \lambda_\gamma + \Gamma_{\beta \alpha}^\gamma (\lambda_\beta - \lambda_\alpha) = 0 \:.
    \end{align}
    Zde $\partial_\alpha$ je derivace ve směru $e_\alpha$ a konexe $\Gamma_{\alpha \beta}^\gamma$ jsou definované rovnostmi
    \begin{align}
        \nabla_{e_\beta} e_\alpha = \sum \Gamma^\gamma_{\alpha \beta} e_\gamma \:.
    \end{align}

    Normalizujme nyní vlastní vektory $e_\alpha$ tak, aby v této bázi byla metrika jednotková matice, tj. $g^{\alpha \beta} = \delta^{\alpha \beta}$.

    Výraz $c^\gamma_{\alpha \beta}$ má význam torze (z definice) a platí
    \begin{align}
        c^\gamma_{\beta \alpha} = \Gamma^\gamma_{\alpha \beta} - \Gamma^\gamma_{\beta \alpha} \:.
    \end{align} 

\end{proof}


\section{Multi-dimensional case}

\begin{proposition}[Changing the space coordinates]
    For any changes of the spatial variables $x^ \alpha \mapsto c^\alpha_\beta x^\beta$, where $\alpha = 1, \cdots, d$ and $\det (c^\alpha_\beta) = 1$, the metrics $g^{ij \alpha}$ and $b^{ij \alpha}_k$ transforms as components of a vector.
\end{proposition}



Ve vícerozměrném případě máme lineární bundle metrik a s nimi spojených konexí. Pro každou záměnu prostorových proměnných $x^\alpha \mapsto c^\alpha_\beta x^\beta$ pro $\alpha = 1, \cdots, d$ splňující $\det c^\alpha_\beta = 1$ se metrika $g^{ij \alpha}$ a konexe $b^{ij \alpha}_k$ transformují jako komponenty vektoru.

\begin{definition}
    Řekneme, že bundle metrik $g^{ij \alpha}$ je silně nedegenerovaný, jestliže pro nějakou sadu $c_\alpha$ je lineární kombinace $c_\alpha g^{ij \alpha}$ nedegenerovaná matice.
\end{definition}

\begin{theorem}
    Nechť je hamiltonovský systém silně nedegenerovaný.
    \begin{enumerate}
        \item Pro $N=1$ lze Poissonovu závorku redukovat na konstantní formu
        \begin{align}
            g^{ij}(u) = \tilde g^{ij}(u) \:.
        \end{align}

        \item Pro $N \geq 2$ lze Poissonovu závorku redukovat na lineární formu
        \begin{align}
            g^{ij \alpha}(u) = g^{ij \alpha}_k u^k + \tilde g^{ij \alpha} \:, \quad \alpha = 1, \cdots , d \:, \\
        \end{align}
        kde koeficienty $g^{ij \alpha}_k = b^{ij \alpha}_k + b^{ij \alpha}_k$, $\tilde g^{ij \alpha}$ a $b^{ij \alpha}_k$ jsou konstantní.
    \end{enumerate}
\end{theorem}

\chapter{Systems of hydrodynamic type:  multi-dimensional case}

When one moves to more dimensions, the theory becomes much more complicated, opposite to the one dimension, where the non-degeneracy of the bracket together with skew-symmetry and Jacobi identity leads not only to the Levi-Civita covariant derivative with zero curvature, but also to the existence of the flat coordinates.
The objectives are the following:
\begin{enumerate}
    \item To find conditions under which the multi-dimensional bracket is Poisson, i.e., satisfies skew-symmetry, Leibniz rule, and Jacobi identity.
    \item To find such coordinates, where the non-degenerate Poisson bracket is reduced to the constant form.
    \item To classify the degenerate brackets.
\end{enumerate}

The first problem was solved in 2006 by O. I. Mochov in \cite{Mochov}, where also the complete classification of nonsingular non-degenerate brackets was found. The theory of compatible metrics plays a key role in the second problem.  The classification of two-dimensional systems was given by Ferapontov, Lorenzoni and Savoldi. The main results are summarised in the doctoral thesis of the latter author \cite{Savoldi}. We shall present some of the ideas here, following mainly the two authors mentioned above.

\section{Multi-dimensional brackets. Mokhov conditions}

In the case with multi-dimensional Dubrovin-Novikov bracket
\begin{align}
    \set{I_1,I_2} = \sum_{\alpha = 1}^ d \int \D x 
    \left[ \frac{\delta I_1}{\delta u^p(x)} \left( g^{pq \alpha}(u) \der{}{x^\alpha} + b^{pq\alpha}_s(u) \pder{u^s}{x^\alpha} \right) \frac{\delta I_2}{\delta u^q(x)} \right] \:, \label{eq:multi-dimensional bracket}
\end{align}
we have a set of metrics $g^{ij \alpha}$ and connections $b^{ij \alpha}_k$, where $\alpha = 1 , \cdots, d$. We can again introduce Christoffel symbols defined by
\begin{align}
    b^{i j \alpha}_k \coloneqq - g^{is \alpha} \Gamma^{j \alpha}_{sk}\:,
\end{align}
(there is no summation over $\alpha$) and the corresponding set of covariant derivatives $\nabla^{\alpha}_j$. The first observation is that these coefficients represent coordinates of a vector under a transformation of the space coordinates $x^\alpha$.
\begin{proposition}[Transformation of the space coordinates]
    For any changes of the spatial variables $\tilde x^\alpha \coloneqq c^\alpha_\beta x^\beta$, where $\alpha = 1, \cdots, d$ and $\det (c^\alpha_\beta) = 1$, the metrics $g^{ij \alpha}$ and $b^{ij \alpha}_k$ transform as components of a vector.
\end{proposition}

\begin{proof}
    The only variable that changes is $\pder{u^i}{x^\alpha}$
    \begin{align}
        \pder{u^i}{x^\beta} = \pder{u^i}{\tilde x^\alpha} \pder{\tilde x^\alpha}{x^\beta} = \pder{u^i}{\tilde x^\alpha} c^\alpha_\beta \:.
    \end{align}
    Hence
    \begin{align}
        g^{ij \alpha}[u(\tilde x)] = c^\alpha_\beta g^{ij \beta}[u(x)] \:, \quad b^{ij \alpha}_k [u(\tilde x)] = c^\alpha_\beta b^{ij \beta}_k [u(x)] \:.
    \end{align}
\end{proof}


The formulation of the generalized form of the Dubrovin-Novikov theorem (Theorem \vref{Theorem-Dubrovin-Novikov-1D}) is given (without proof) by Mokhov in \cite{Mochov}.

\begin{theorem}[Mokhov conditions; \cite{Mochov}]
    A Dubrovin-Novikov bracket \eqref{eq:Dubrovin-Novikov bracket} is a Poisson bracket if and only if the following relations for the coefficients are fulfilled:
    \begin{align}
        g^{ij \alpha} = g^{ji \alpha} \:, \label{eq:Mochov-A} \\
        g^{ij \alpha}_{,k} = b^{ij \alpha}_k + b^{ji \alpha}_k \:, \label{eq:Mochov-B}
    \end{align}
    \begin{align}
        \sum_{(\alpha, \beta)} (g^{mi \alpha} b^{jk \beta}_m - g^{m j \beta} b^{i k \alpha}_m) = 0 \:, \label{eq:Mochov-C}\\
        \sum_{(i,j,k)} (g^{mi \alpha} b^{jk \beta}_m - g^{m j \beta} b^{ik \alpha}_m) = 0 \:, \label{eq:Mochov-D}\\
        \sum_{(\alpha, \beta)} \left[ g^{si \alpha} (b^{jk \beta}_{s,q} - b^{jk \beta}_{q,s} ) + b^{ij \alpha}_s b^{s k \beta}_q - b^{i k \alpha}_s b^{k j \beta}_q \right] = 0 \label{eq:Mochov-E}\:,
    \end{align}
    \begin{align}
        g^{si \beta} b^{j k \alpha}_{q,s} - b^{ij \beta}_s b^{s k \alpha}_q - b^{i k \beta}_s b^{j s \alpha}_q = g^{s j \alpha} b^{i k \beta}_{q,s} - b^{ji \alpha}_s b^{s k \beta}_q - b^{is \beta}_q b^{jk \alpha}_s \label{eq:Mochov-F}\:,
    \end{align}
    \begin{align}
        \pder{}{u^r} \left[ g^{si \alpha} (b^{j k \beta}_{s,q} - b^{j k \beta}_{q,s}) + b^{ij \alpha}_s b^{s k \beta}_q - b^{i k \alpha}_s b^{s j \beta}_q \right] 
        \nonumber \\+ \pder{}{u^q} \left[ g^{si \beta} (b^{j k  \alpha}_{s,k} - b^{j k \alpha}_{r,s}) + b^{ij \beta}_s b^{s k \alpha}_{s,r} + b^{ij \beta}_s b^{s k \alpha}_r - b^{i k \beta}_s b^{k j \alpha}_s \right] \label{eq:Mochov-G}
        \\+ \sum_{(i,j,k)} \left[ b^{si \beta}_q (b^{j k \alpha}_{r,s} - b^{j k \alpha}_{s,r}) + b^{s i \alpha}_r (b^{j k \beta}_{q,s} - b^{j k \beta}_{s,q}) \right] \nonumber
        = 0 \:.
    \end{align}
\end{theorem}
The summations $\sum_{(\alpha, \beta)}$ and $\sum_{(i,j,k)}$ mean summations over all cyclic permutations of the indicated indices. Note that in \eqref{eq:Mochov-C}-\eqref{eq:Mochov-G} the indices $\alpha$ and $\beta$ do not have to be distinct indices (the cyclic permutation $\sum_{(\alpha, \beta)}$ can be ommited in that case). Thus, in particular, we get an important lemma.
\begin{lemma}
    Every summand on the right-side of the multi-dimensional bracket \eqref{eq:multi-dimensional bracket} is one-dimensional Dubrovin-Novikov bracket \eqref{eq:Dubrovin-Novikov bracket 1D}, i.e., each multidimensional bracket of the form \eqref{eq:multi-dimensional bracket} is always the sum of one-dimensional brackets \eqref{eq:Dubrovin-Novikov bracket 1D} with respect to each of the independent variables $x^\alpha$.
\end{lemma}


\begin{remark}
    Let us take once again a closer look to the case of $d=1$ and compare it with the Grinberg conditions (\vref{Grinberg conditions}). The Mokhov conditions are following:
    \begin{align}
        g^{ij} =& g^{ji} \quad &\textbf{(symmetry of metric)} \:, \label{eq:Mochov1D-A}\\
        g^{ij}_{,k} =& b^{ij}_k + b^{ji}_k \quad &\textbf{(compatibility condition)} \label{eq:Mochov1D-B}\:, \\
        g^{mi} b^{jk}_m - g^{mj} b^{ik}_m =& 0 \quad &\textbf{(zero torsion)} \:, \label{eq:Mochov1D-C}\\
        \sum_{(i,j,k)} g^{mi} b^{jk}_m - g^{mj} b^{ik}_m =& 0 \:,\quad& \label{eq:Mochov1D-D}
    \end{align}
    \begin{align}
        g^{si} (b^{jk}_{s,q} - b^{jk}_{q,s}) + b^{ij}_s b^{sk}_q - b^{ik}_s b^{kj}_q = 0 \quad &\textbf{(zero curvature)}\:, \label{eq:Mochov1D-E}
    \end{align}
    \begin{align}
        g^{si} b^{j k}_{q,s} - b^{ij}_s b^{s k}_q - b^{i k}_s b^{j s}_q = g^{s j } b^{i k}_{q,s} - b^{ji}_s b^{s k}_q - b^{is}_q b^{jk}_s \:, \label{eq:Mochov1D-F}\\
        \pder{}{u^r} \left[ g^{si} (b^{j k }_{s,q} - b^{j k}_{q,s}) + b^{ij}_s b^{s k}_q - b^{i k}_s b^{s j}_q \right] 
        \nonumber \\+ \pder{}{u^q} \left[ g^{si } (b^{j k}_{s,k} - b^{j k}_{r,s}) + b^{ij \beta}_s b^{s k}_{s,r} + b^{ij}_s b^{s k}_r - b^{ik}_s b^{kj}_s \right] \label{eq:Mochov1D-G}
        \\+ \sum_{(i,j,k)} \left[ b^{si}_q (b^{j k}_{r,s} - b^{j k}_{s,r}) + b^{s i}_r (b^{j k}_{q,s} - b^{j k}_{s,q}) \right] \nonumber
        = 0 \:.
    \end{align}
    It is evident that \eqref{eq:Mochov1D-D} is just \eqref{eq:Mochov1D-C} with a cyclic permutation, hence it holds trivially. It is not so obvious, but also true, that \eqref{eq:Mochov1D-F} follows from the compatibility condition, zero torsion and zero curvature. 
     The last relation \eqref{eq:Mochov1D-G} is the one remaining Grinberg condition \eqref{eq:Grinberg-5}. Only if we assume non-degeneracy, it follows from zero torsion and curvature, otherwise it must be involved as a separate assumption.
\end{remark}


\section{Tensor obstructions}

\begin{definition}[The obstruction tensors]
    We define the obstruction tensors
    \begin{align}
        T^{i \alpha \beta}_{jk} (u) \coloneqq \Gamma^{i \beta}_{jk}(u) - \Gamma^{i \alpha}_{jk}(u) \:,
    \end{align}
    for each pair of distinct indices $\alpha$ and $\beta$.
\end{definition}
We shall also use their fully contravariant form
\begin{align}
    T^{ijk \alpha \beta}(u) \coloneqq g^{i r \alpha}(u) g^{ks \beta}(u) T^{j \alpha \beta}_{rs} (u) \:.
\end{align}
Note that $T^{i \alpha \beta}_{jk} (u)$ are transformed as tensors only in the upper $\alpha$, $\beta$ indices (since in the remaining they transform as Christoffel symbols).

\begin{lemma} \label{lemma:T=0}
    Non-degenerate Dubrovin-Novikov brackets can be reduced to a constant bracket by a local change of coordinates if and only if all the obstruction tensors $T^{ij \alpha}_{jk}(u)$ are identically equal to zero.
\end{lemma}
\begin{proof}
    For any constant bracket all the connections $\Gamma^{i\alpha}_{jk}(u)$ are equal to zero identically, therefore all obstruction tensors are zero. Conversly, if all the $T^{ij \alpha}_{jk}(u)$ are identically zero, then all the connections $\Gamma^{i\alpha}_{jk}(u)$ are equal to each other. Consequently, they are all zero in any flat coordinates of the metric $g^{ij 1}(u)$. All the remaining metrics $g^{ij \beta}(u)$ are necessarily constant in these coordinates by virtue of compatiblity of the metrics with the corresponding connections.
\end{proof}


\begin{theorem}[Mokhov; \cite{Mochov}]
    Flat non-degenerate metrics $g^{ij \alpha}$ define a multidimensional Poisson bracket if and only if the following relations are fulfilled:
    \begin{align}
        T^{ijk \alpha \beta}(u) =& T^{kji \alpha \beta}(u) \:, \label{eq:Mochov-1}\\
        \sum_{(i,j,k)} T^{ijk \alpha \beta}(u) =& 0 \:, \label{eq:Mochov-2}\\
        T^{ij s \alpha \beta}(u) T^{k \alpha \beta}_{st}(u) =& T^{i k s \alpha \beta}(u) T^{j \alpha \beta}_{st} (u) \:, \label{eq:Mochov-3}\\
        \nabla^\alpha_r T^{ijk \alpha \beta}(u) =& 0 \:, \label{eq:Mochov-4}
    \end{align}
    where $\nabla^\alpha_k$ is the Levi-Civita covariant derivative generated by the metric $g^{ij \alpha}(u)$ and $\Gamma^{i \alpha}_{jk}(u)$ are the corresponding Christoffel symbols. 
\end{theorem}

\begin{remark}
    The relations \eqref{eq:Mochov-1} and \eqref{eq:Mochov-3} were found by Dubrovin and Novikov in \cite{Dubrovin-Novikov}. In \cite{Mochov} it was proved that this set of tensor relations is not complete and the remaining two were obtained.
\end{remark}

\begin{proof}
    The first equation \eqref{eq:Mochov-1} reads as
    \begin{align}
        g^{ir \alpha} g^{ks \beta} (\Gamma^{j\beta}_{rs} - \Gamma^{j \alpha}_{rs}) = g^{kr \alpha} g^{is \beta} (\Gamma^{j\beta}_{rs} - \Gamma^{j \alpha}_{rs}) \:,
    \end{align}
    which is equivalent to
    \begin{align}
        \Gamma^{j \beta}_{rs} - \Gamma^{j \alpha}_{rs} = 0 \:,
    \end{align}
    and it follows from \eqref{eq:Mochov-C}.

    The equation \eqref{eq:Mochov-2} is equivalent to \eqref{eq:Mochov-D}, \eqref{eq:Mochov-3} is equivalent to \eqref{eq:Mochov-E}, \eqref{eq:Mochov-4} is equivalent to \eqref{eq:Mochov-F}. The last condition \eqref{eq:Mochov-G} is a direct consequence of the relations \eqref{eq:Mochov-A}-\eqref{eq:Mochov-F} in the non-degenerate case.
\end{proof}

\begin{example}[Classification of one-component brackets; \cite{Mochov}]

It is not complicated to obtain a complete classification of all one-component brackets
\begin{align}
    \pder{u}{t} = \left( g^{\alpha}(u) \pder{^2h}{u^2} + b^{\alpha}(u) \pder{h}{u} \right) \pder{u}{x^\alpha} \:, \quad \alpha = 1, \cdots, d \:.
\end{align}
The multi-dimensional bracket is 
\begin{align}
    \set{u(x),u(y)} = g^{\alpha}(u) \partial_\alpha \delta(x-y) + b^\alpha(u) \delta(x-y) \:.
\end{align}
and the decomposition is
\begin{align}
    b^\alpha = - g^\alpha \Gamma^\alpha \:.
\end{align}
The Mokhov conditions \eqref{eq:Mochov-A}, \eqref{eq:Mochov-C}, \eqref{eq:Mochov-E} and \eqref{eq:Mochov-G} are automatically fulfilled. The relations \eqref{eq:Mochov-B} and \eqref{eq:Mochov-D} are
\begin{align}
    \pder{g^\alpha}{u} = 2 b^\alpha(u) \:, \label{eq:1comp-1} \\
    g^\alpha(u) b^\beta(u) = g^\beta(u) g^\alpha(u) \label{eq:1comp-2}\:.
\end{align}
The last condition \eqref{eq:Mochov-F} follows from them. If we substitute \eqref{eq:1comp-1} into \eqref{eq:1comp-2}, we get
\begin{align}
    \frac{1}{g^\alpha} \pder{g^\alpha}{u}= \frac{1}{g^\beta}\pder{g^\beta}{u} \:,
\end{align}
hence
\begin{align}
    g^\alpha(u) = 2c^\alpha g(u) \:, \quad b^\alpha(u) = c^\alpha \quad \forall \alpha = 1, \cdots d \:,
\end{align}
where $c^\alpha$ are arbitrary constants.
Therefore, for $g(u) \neq 0$ all the obstruction tensors are identically zero:
\begin{align}
    T^{\alpha \beta} = \Gamma^\beta - \Gamma^\alpha = -\frac{b^\alpha}{g^\alpha} + \frac{b^\beta}{g^\beta} = - \frac{1}{2 g} + \frac{1}{2g} = 0 \:.
\end{align}
By the virtue of lemma \vref{lemma:T=0} by local changing of coordinates we can always find a transform in witch the bracket has a constant form. The 
\begin{align}
    \pder{w}{t} = \eta^\alpha \pder{w}{x^\alpha} \:.
\end{align}

The degenerate case is also trivial. If for some $\alpha$ is $g^\alpha(u) = 0$, then from \eqref{eq:Mochov-B} follows $b^\alpha(u) = 0$.

\end{example}

\section{Compatible brackets}

In this section, we introduce the theory of compatible metrics constructed by Mochov and show that the study of non-degenerate multi-dimensional brackets is connected to the study of compatible one-dimensional brackets. All the proofs can be found in \cite{Mochov}. 

\begin{definition}[Compatible metrics]
    Let $g^{ij1}(u)$ and $g^{ij2}(u)$ be non-degenerate psudo-Riemannian metrics with Levi-Civita connections $\Gamma^{ij1}_k(u)$, $\Gamma^{ij2}_k(u)$. We say that $g^{ij1}$ and $g^{ij2}$ are compatible if for any linear combination of these metrics
    \begin{align}
        g^{ij}(u) = \lambda_1 g^{ij1}(u) + \lambda_2 g^{ij2}(u) \:,
    \end{align}
    where $\lambda_{1,2}$ are arbitrary constants such that $\det g^{ij} \neq 0$, the coefficients of the corresponding Levi-Civita connections and the components of the corresponding Riemann tensors are related by the same linear formula:
    \begin{align}
        \Gamma^{ij}_k = \lambda_1 \Gamma^{ij1}_k + \lambda_2 \Gamma^{ij2}_k \:, \label{eq:compatible-Gamma}\\
        R\indices{_a_b^k_c} = \lambda_1 R\indices{_a_b^k_c^1} + \lambda_2 R\indices{_a_b^k_c^2} \label{eq:compatible-R}\:.
    \end{align}
    If for any linear combination of metrics only \eqref{eq:compatible-Gamma} is fulfilled, then the metrics are called almost compatible.
\end{definition}


\begin{definition}[Compatible Poisson brackets]
    We say that two Poisson brackets are compatible if each their linear combination is a Poisson bracket.
\end{definition}

\begin{definition}[L-affinor]
    Let $g^{ij1}(u)$ and $g^{ij2}(u)$ be two metrics.
    We define the L-affinor $L^i_j$ by
    \begin{align}
        L^i_j(u) = g^{ik2}(u) g_{kj1}(u) \:,
    \end{align}
\end{definition}

\begin{theorem}[Vanishing of the Nijenhuis tensor]
    The metrics $g^{ij1}(u)$ and $g^{ij2}(u)$ are almost compatible if and only if the Nijenhuis tensor of the L-affinor is zero, i.e.,
    \begin{align}
        N^i_{jk} \coloneqq L^s_j L^{i}_{k,s} - L^s_k L^{i}_{j,s} - L^i_s L^s_{k,j} - L^i_s L^s_{j,k} = 0 \:.
    \end{align}
\end{theorem}

Note that there are examples of metrics that are only almost compatible, but not compatible, and for which the Nijenhuis tensor is zero \cite{Mochov-compatibility}]. 

\begin{theorem}[Mutual compatibility of the metrics; \cite{Mochov}] 
    All metrics $g^{ij \alpha}(u)$, $1 \leq \alpha \leq N$, defining a multi-dimensional Poisson bracket \eqref{eq:multi-dimensional bracket} are mutually compatible. All one-dimensional Dubrovin-Novikov brackets forming a multi-dimensional bracket are mutually compatible.
\end{theorem}

\begin{definition}[Nonsingular pair of metrics]
    A pair of metrics $g^{ij 1}(u)$ and $g^{ij 2}(u)$ is called nonsingular if all eigenvalues of this pair of metrics, i.e., the roots of the equation $\det (g^{ij 1}(u)- \lambda g^{ij 2}(u)) = 0$, are distinct.
\end{definition}

\begin{theorem}[Mokhov classification theorem; \cite{Mochov}]
    If for a non-degenerate multidimensional Poisson bracket one of the metrics $g^{ij \alpha}(u)$ forms nonsingular pairs with all the remaining metrics, then this Poisson bracket can be reduced to a constant form by a local change of coordinates.
\end{theorem}

The theory of singular non-degenerate brackets  and also all degenerate brackets is still not complete at present. 
%\chapter*{Conclusion}
\addcontentsline{toc}{chapter}{Conclusion}

The area of geometry in hydrodynamics offers many challenging, unsolved problems. This thesis aimed to explore the Dubrovin-Novikov theorem which shows that for non-degenerate metric the existence of Poisson structure is equivalent to the existence of Levi-Civita connection with zero curvature, therefore connects the Poisson and pseudo-Riemannian geometry. The proof of this theorem in the original paper \cite{Dubrovin-Novikov} was only sketched, so the potential benefit of this thesis was to provide it in perhaps a readable, detailed form. In Chapter 3, we have first verified the tensor structure of the functions in quasi-linear systems of the first order and constructed the Dubrovin-Novikov bracket for systems of hydrodynamic type. Then we restricted ourselves to the one-dimensional case, where we have shown that the functions forming the non-degenerate bracket are in fact inverse metric and Christoffel symbols. Then the proof of the Dubrovin-Novikov theorem was divided into four steps, in which we have 1. connected the skew-symmetry of the bracket with the symmetry of the metric and the compatibility condition, 2. written down the conditions under which the Jacobi identity holds, 3. verified the zero torsion and curvature and 4. indicated the existence of the coordinates in which the bracket is trivial and therefore finished the converse statement. In Chapter 4, we have done a similar analysis for the multi-dimensional bracket, written down the conditions under which the bracket is Poisson, and reformulated the problem of finding the flat coordinates via the obstruction tensors. We end the discussion with the theorem by Mokhov, which states that the existence of flat coordinates is possible if one of the metrics forms nonsingular pairs with all the remaining metrics. We did not explore whether this condition is satisfied in the examples mentioned in Chapter 2.


Nevertheless, the connection between the two geometric structures is still slightly mysterious, as the zero torsion and zero curvature emerge among a huge amount of calculations. Is there something deeper to be identified beyond these equations? 

The assumption which is found to be crucial is the non-degeneracy of a metric. We are convinced that there exist attempts to describe the geometry with degenerate metrics. The problem with hydrodynamical systems is that the functions $g^{ij \alpha}$ are in fact the inverse metrics; therefore we need a theory which stems from the symmetric but degenerate (2,0) tensor. If the metric is degenerate, there is no Levi-Civita connection. There could be several covariant derivatives that annihilate the metric and have zero torsion. Could this be the way of proving the symmetric hyperbolicity for the degenerate systems? At least to the author's knowledge, these questions are still open to inquiry. 

To the inquisitive reader, we further recommend two sources that widen the horizons of current knowledge. For an analysis of two-dimensional systems, see \cite{Savoldi}, and for more of the one-dimensional degenerate brackets, see \cite{savoldi2014deformations}. 


%%% Bibliography
%%% Bibliography (literature used as a source)
%%%
%%% We employ bibTeX to construct the bibliography. It processes
%%% citations in the text (e.g., the \cite{...} macro) and looks up
%%% relevant entries in the bibliography.bib file.
%%%
%%% The \bibliographystyle command selects, which style will be used
%%% for references from the text. The argument in curly brackets is
%%% the name of the corresponding style file (*.bst). Both styles
%%% mentioned in this template are included in LaTeX distributions.

%\bibliographystyle{plainnat}    %% Author (year)
\bibliographystyle{unsrt}     %% [number]
\bibliography{bibliography}
\renewcommand{\bibname}{Bibliography}

%%% Generate the bibliography. Beware that if you cited no works,
%%% the empty list will be omitted completely.

%%% If case you prefer to write the bibliography manually (without bibTeX),
%%% you can use the following. Please follow the ISO 690 standard and
%%% citation conventions of your field of research.

% \begin{thebibliography}{99}
%
% \bibitem{lamport94}
%   {\sc Lamport,} Leslie.
%   \emph{\LaTeX: A Document Preparation System}.
%   2nd edition.
%   Massachusetts: Addison Wesley, 1994.
%   ISBN 0-201-52983-1.
%
% \end{thebibliography}

%%% Figures used in the thesis (consider if this is needed)
%\listoffigures

%%% Tables used in the thesis (consider if this is needed)
%%% In mathematical theses, it could be better to move the list of tables to the beginning of the thesis.
% \listoftables

%%% Abbreviations used in the thesis, if any, including their explanation
%%% In mathematical theses, it could be better to move the list of abbreviations to the beginning of the thesis.
% \chapwithtoc{List of Abbreviations}

%%% Attachments to the bachelor thesis, if any. Each attachment must be
%%% referred to at least once from the text of the thesis. Attachments
%%% are numbered.
%%%
%%% The printed version should preferably contain attachments, which can be
%%% read (additional tables and charts, supplementary text, examples of
%%% program output, etc.). The electronic version is more suited for attachments
%%% which will likely be used in an electronic form rather than read (program
%%% source code, data files, interactive charts, etc.). Electronic attachments
%%% should be uploaded to SIS and optionally also included in the thesis on a~CD/DVD.
%%% Allowed file formats are specified in provision of the rector no. 72/2017.
%\appendix
%\chapter{Appendix}

%\openright
\end{document}
